\documentclass[11pt]{article}
\usepackage[a4paper,margin=28mm]{geometry}
\usepackage{amsmath,amssymb,amsthm}
\usepackage[colorlinks=true,linkcolor=blue,urlcolor=blue,citecolor=blue]{hyperref}
\newtheorem*{theorem*}{Theorem}
\newtheorem*{corollary*}{Corollary}
\newenvironment{statement}[1]{\par\medskip\noindent\phantomsection\textbf{#1}\ }{\par\medskip}
\setlength{\emergencystretch}{3em}
\newcommand{\Q}{\mathbb Q}
\newcommand{\Z}{\mathbb Z}
\newcommand{\F}{\mathbb F}
\DeclareMathOperator{\End}{End}
\DeclareMathOperator{\Gal}{Gal}
\DeclareMathOperator{\GL}{GL}
\DeclareMathOperator{\PGL}{PGL}
\DeclareMathOperator{\Hom}{Hom}
\DeclareMathOperator{\NS}{NS}
\DeclareMathOperator{\Sym}{Sym}
\DeclareMathOperator{\tr}{tr}
\title{No rational \(13\)-torsion on geometrically split Jacobians that are simple over \(\Q\)}
\author{}
\date{21 October 2026}

\begin{document}
\maketitle

\begin{theorem*}
Let \(C/\Q\) be a curve of genus \(2\) whose Jacobian \(J\) is geometrically split. If
      \(J(\Q)\) contains a point of order \(13\), then \(J\) is not simple over \(\Q\).
\end{theorem*}

\begin{corollary*}
No geometrically split genus-\(2\) Jacobian over \(\Q\) has a rational point of
      order \(13\). In particular, \href{https://genus-2-torsion.github.io/group.html?g=13}{\([13]\)}, \href{https://genus-2-torsion.github.io/group.html?g=26}{\([26]\)}, \href{https://genus-2-torsion.github.io/group.html?g=2.26}{\([2,26]\)} and \href{https://genus-2-torsion.github.io/group.html?g=39}{\([39]\)} do not occur as the rational torsion
      subgroup of a geometrically split genus-\(2\) Jacobian over \(\Q\).
\end{corollary*}

Here ``geometrically split'' means that \(J_{\bar\Q}\) is isogenous to a product of two elliptic curves;
    we impose no restriction on the degree of this isogeny, on the field of definition of the geometric
    endomorphisms, or on the integral endomorphism ring. The proof uses the principal polarization of \(J\),
    and this cannot be avoided: there are abelian surfaces over \(\Q\) that are simple over \(\Q\), geometrically
    split, and have a rational point of order \(13\) (Section~\ref{remarks}). Proposition 2 shows that each of them is of one of two special types.

The Magma scripts and logs cited below are published with this proof at \url{https://genus-2-torsion.github.io/data/knowledge/sources/}.

\section{Conventions and standard facts}\label{facts}

For a field \(k\) we write \(G_k=\Gal(\bar k/k)\). For a prime \(\ell\), \(\chi_\ell\) is the
      \(\ell\)-adic cyclotomic character and \(\bar\chi_\ell\) its reduction modulo \(\ell\); thus
      \(\bar\chi_{13}:G_{\Q}\to\F_{13}^\times\) is surjective of order \(12\). For an abelian variety \(X\)
      over a field \(k\) of characteristic \(0\), \(T_\ell X\) is its \(\ell\)-adic Tate module,
      \(V_\ell X=T_\ell X\otimes\Q_\ell\), \(\End^0=\End\otimes\Q\), \(\Hom^0=\Hom\otimes\Q\), and
      \(X\sim_k Y\) means that \(X\) and \(Y\) are isogenous over \(k\). For a \(G\)-stable
      \(\Z_\ell\)-lattice \(\Lambda\) in a \(\Q_\ell[G]\)-module we put \(\bar\Lambda=\Lambda/\ell\Lambda\),
      and \((\cdot)^{\mathrm{ss}}\) denotes semisimplification. A ``constituent'' is a Jordan--Hölder factor.
      For a character \(\eta\) of \(G_K\) and \(\tau\in G_{\Q}\) normalizing \(G_K\) we put
      \(\eta^\tau(g)=\eta(\tau^{-1}g\tau)\), and similarly for representations.

We use the following standard facts.

\begin{itemize}
\item \textbf{(F1) Lattices and isogenies.} Let \(A\) be an abelian variety over a number field \(k\).
          A \(k\)-isogeny \(f:B\to A\) gives \(G_k\)-stable lattices \(f_*T_\ell B\subset V_\ell A\), equal to
          \(T_\ell A\) for almost all \(\ell\). If \(f:B\to A\) and \(f':B'\to A\) are \(k\)-isogenies with
          \(f'_*T_\ell B'=\alpha\, f_*T_\ell B\) for every \(\ell\), for some \(\alpha\in\End^0_k(A)^\times\),
          then \(B\cong_k B'\): the quasi-isogeny \(f^{-1}\alpha^{-1}f'\) is an isomorphism. For an
          elliptic curve \(E/k\), every \(G_k\)-stable lattice \(\Lambda\subset V_\ell E\) is \(h_*^{-1}(T_\ell P)\)
          for some \(k\)-isogeny \(h:E\to P\) of \(\ell\)-power degree, up to multiplication by a power of \(\ell\).
\item \textbf{(F2) Riemann forms.} A polarization \(\lambda\) of an abelian variety \(B/\Q\) gives a
          \(G_{\Q}\)-equivariant alternating pairing \(E^\lambda:T_\ell B\times T_\ell B\to\Z_\ell(1)\),
          nondegenerate on \(V_\ell B\). The polarization is principal if and only if \(E^\lambda\) is perfect
          on \(T_\ell B\) for every \(\ell\). The class of \(\lambda\) lies in
          \(\NS(B_{\bar\Q})^{G_{\Q}}\), and a quasi-isogeny identifies the spaces
          \(\NS(\cdot_{\bar\Q})\otimes\Q\) compatibly with the Galois action and the Riemann forms.
\item \textbf{(F3) Néron--Severi group of a square.} Let \(E\) be an elliptic curve without CM over an
          algebraically closed field, with Weil pairing \(e\). Then \(\NS(E\times E)\cong\Sym_2(\Z)\), the
          symmetric integral \(2\times2\) matrices: \(S\) corresponds to the class whose Riemann form on
          \(T_\ell E\otimes\Z^2\) is \((v\otimes a,w\otimes b)\mapsto a^{t}Sb\cdot e(v,w)\). A positive definite
          \(S\) is a polarization of degree \((\det S)^2\), and \(S=I\) is the product polarization.
          An endomorphism \(N\in M_2(\Z)=\End(E\times E)\) has degree \((\det N)^2\).
\item \textbf{(F4) Brauer--Nesbitt.} For a continuous representation of a profinite group \(G\) on a
          finite-dimensional \(\Q_\ell\)-vector space \(V\), the semisimplification \(\bar\Lambda^{\mathrm{ss}}\)
          does not depend on the \(G\)-stable lattice \(\Lambda\subset V\), and the constituents of any
          quotient \(\Lambda_1/\Lambda_2\) of stable lattices \(\Lambda_2\subset\Lambda_1\) lie among those of
          \(\bar\Lambda^{\mathrm{ss}}\). The same applies to two isomorphic representations, after
          transporting a lattice through the isomorphism.
\item \textbf{(F5) Burnside and Nakayama.} If a subgroup \(H\subset\GL_n(\Z_\ell)\) has absolutely
          irreducible reduction modulo \(\ell\), then the \(\Z_\ell\)-span of \(H\) is \(M_n(\Z_\ell)\):
          by Burnside's theorem its reduction spans \(M_n(\F_\ell)\), and Nakayama's lemma lifts this.
\item \textbf{(F6) Morita.} Let \(\mathcal O\) be \(\Z_\ell\), a finite étale \(\Z_\ell\)-algebra, or
          the ring of integers of a finite extension of \(\Q_\ell\), and let \(W=\mathcal O^2\) with its
          natural \(M_2(\mathcal O)\)-action. Every \(M_2(\mathcal O)\)-stable \(\Z_\ell\)-lattice
          \(\Lambda\subset W\otimes\Q\) equals \(\mathfrak aW\) for a fractional \(\mathcal O\)-ideal
          \(\mathfrak a\). Indeed \(\Lambda=e_{11}\Lambda\oplus e_{22}\Lambda\), the matrix unit \(e_{21}\)
          maps \(e_{11}\Lambda\) isomorphically onto \(e_{22}\Lambda\), and \(e_{11}\Lambda\) is an
          \(\mathcal O\)-lattice in a free module of rank one over \(\mathcal O\otimes\Q\), hence of the form
          \(\mathfrak a\,e_{11}W\).
\item \textbf{(F7) Forms of a square.} Let \(E\) be an elliptic curve without CM over a number field
          \(k\). Then \(\End^0(E^2_{\bar k})=M_2(\Q)\) consists of endomorphisms defined over \(k\), and
          abelian varieties \(X/k\) with \(X_{\bar k}\sim E^2_{\bar k}\) correspond, up to \(k\)-isogeny, to
          \(\GL_2(\Q)\)-conjugacy classes of continuous homomorphisms \(\rho:G_k\to\GL_2(\Q)\), which have
          finite image. We write \(X\sim_k E\otimes\rho\). Then \(G_k\) acts on
          \(\End^0(X_{\bar k})\cong M_2(\Q)\) by \(\mathrm{Ad}\circ\rho\), and \(\End^0_k(X)\) is the
          commutant of \(\rho(G_k)\) in \(M_2(\Q)\). If \(\Lambda\cong\Z^2\) is \(\Z[i]\) or \(\Z[\zeta_3]\)
          and \(\psi:G_k\to\Lambda^\times\) is a character of finite order, then \(E\otimes\Lambda(\psi)\)
          denotes the twist of \(E\otimes_{\Z}\Lambda\cong E^2\) by the cocycle \(g\mapsto\psi(g)\) (acting
          by multiplication on \(\Lambda\)); its Tate module is \(T_\ell E\otimes\Lambda\), on which \(g\)
          acts by \(\rho_{E,\ell}(g)\otimes\psi(g)\).
\item \textbf{(F8) Product polarizations are not Jacobians.} If \((B,\lambda)\) is a principally polarized
          abelian surface over \(\Q\) such that \((B_{\bar\Q},\lambda)\) is isomorphic to a product
          \(E_1\times E_2\) with the product polarization, then \((B,\lambda)\) is not the Jacobian of a
          curve of genus \(2\): the theta divisor of a genus-\(2\) Jacobian is irreducible, while that of
          a product is \(E_1\times\{0\}\cup\{0\}\times E_2\).
\item \textbf{(F9) Elliptic curves over \(\Q\).} No elliptic curve over \(\Q\) has a rational point of
          order \(13\) (Mazur~\cite[Theorem (7\({}'\))]{Maz}). No elliptic curve over \(\Q\) has a rational cyclic
          isogeny of degree \(26\), \(39\), or \(169\): for \(26\) and \(39\) this is part of the computations
          \href{https://genus-2-torsion.github.io/data/knowledge/sources/qsimple_13_x0_26.m}{\texttt{qsimple\_\allowbreak{}13\_\allowbreak{}x0\_\allowbreak{}26.m}} and \href{https://genus-2-torsion.github.io/data/knowledge/sources/qsimple_13_x0_39.m}{\texttt{qsimple\_\allowbreak{}13\_\allowbreak{}x0\_\allowbreak{}39.m}} (the rational points of \(X_0(26)\) and
          \(X_0(39)\) are cusps; for \(39\) see also Kenku~\cite{Ken79}), and for \(169\) it is
          Kenku~\cite{Ken80}. In particular, an elliptic curve over \(\Q\) has at most one rational
          subgroup of order \(13\): two of them, \(C\neq C'\), would give the cyclic isogeny
          \(E/C'\to E\to E/C\) of degree \(169\).
\item \textbf{(F10) Faltings.} For abelian varieties \(X,Y\) over a number field \(k\), \(V_\ell X\) is a
          semisimple \(G_k\)-module and \(\Hom_k(X,Y)\otimes\Z_\ell\to{}\) \(\Hom_{G_k}(T_\ell X,T_\ell Y)\) is an
          isomorphism (Faltings~\cite[§5, Satz 3 and Korollar 1]{Fal}). In particular, for an elliptic
          curve \(E/k\) without CM, \(V_\ell E\) is an irreducible \(G_k\)-module.
\end{itemize}

\begin{statement}{Lemma 1.}\label{lemma-sequence}
Let \(X\) be an abelian surface over a number field \(k\), let \(E\subset X\) be an elliptic curve
          defined over \(k\), and let \(Q\in X(k)\) have prime order \(q\). Then either \(Q\in E(k)\), or the
          image of \(Q\) in \((X/E)(k)\) has order \(q\).
\end{statement}

\emph{Proof.} The image of \(Q\) is zero exactly when \(Q\in E\). \qed

\section{The case distinction and the main proposition}\label{cases}

Let \(A/\Q\) be an abelian surface that is simple over \(\Q\), with
      \(A_{\bar\Q}\sim E_1\times E_2\). Exactly one of the following holds:

\begin{itemize}
\item \textbf{(N)} \(E_1\) and \(E_2\) are not isogenous;
\item \textbf{(CM)} \(E_1\sim E_2=:E\) and \(E\) has complex multiplication;
\item \textbf{(S)} \(E_1\sim E_2=:E\) and \(E\) has no complex multiplication.
\end{itemize}

In case (S), \(D:=\End^0(A_{\bar\Q})\cong M_2(\Q)\). Let \(L\) be the smallest field over which all
      geometric endomorphisms of \(A\) are defined; it is a finite Galois extension of \(\Q\). The group
      \(\Gamma=\Gal(L/\Q)\) acts faithfully on \(D\), and \(\End^0_{\Q}(A)=D^\Gamma\). By Skolem--Noether
      this action embeds \(\Gamma\) into \(\operatorname{Aut}(M_2(\Q))=\PGL_2(\Q)\); we write
      \(g\mapsto[R_g]\), so that \(g\) acts on \(D\) by \(\mathrm{Ad}(R_g)\).

A finite subgroup of \(\PGL_2(\mathbb R)\) is cyclic or dihedral: it fixes a point of the
      hyperbolic plane and embeds into \(O(2)\). If \([R]\in\PGL_2(\Q)\) has order \(n>2\), the ratio of
      the eigenvalues of \(R\) is a primitive \(n\)-th root of unity \(\zeta\), and
      \((\tr R)^2/\det R=2+\zeta+\zeta^{-1}\in\Q\); hence \(n\in\{3,4,6\}\). Therefore

\[
\Gamma\in\{C_1,\ C_2,\ C_3,\ C_4,\ C_6,\ V_4,\ D_6,\ D_8,\ D_{12}\},
\]

where \(D_{2n}\) is dihedral of order \(2n\) (so \(D_6\cong S_3\)) and \(V_4\) is the Klein four group.
      An element \(r\) of order \(n\in\{3,4,6\}\) has a unique representative \(R\) with \(\tr R=1\); its
      characteristic polynomial is \(X^2-X+1/c\) with \(c=1,2,3\) respectively, of discriminant
      \(-3,-1,-1/3\).

Every elliptic subvariety of \(A_{\bar\Q}\) is the image of an integral multiple of a rank-one
      idempotent of \(D\), so it is defined over \(L\). If \(E,E'\subset A_L\) are elliptic subvarieties,
      every geometric homomorphism \(E'\to E\) is defined over \(L\): composing with an \(L\)-rational
      quasi-projection \(A\to E'\) and the inclusion \(E\subset A\) embeds \(\Hom^0(E'_{\bar\Q},E_{\bar\Q})\)
      \(G_L\)-equivariantly into \(D\), on which \(G_L\) acts trivially.

\begin{statement}{Proposition 2.}\label{prop-main}
Let \(A/\Q\) be an abelian surface that is simple over \(\Q\) and geometrically split, with
          \(A(\Q)[13]\neq0\). Then \(A\) is of type (S) and \(\Gamma\cong C_2\) or \(\Gamma\cong C_3\). Moreover:

\begin{enumerate}
\renewcommand{\labelenumi}{(\alph{enumi})}
\item if \(\Gamma\cong C_3\), no abelian surface that is \(\Q\)-isogenous to \(A\) admits a principal
              polarization over \(\Q\);
\item if \(\Gamma\cong C_2\), then for every principally polarized abelian surface \((B,\lambda)\) over
              \(\Q\) with \(B\sim_{\Q}A\), the pair \((B_{\bar\Q},\lambda)\) is isomorphic to a product of two
              elliptic curves with the product polarization.
\end{enumerate}
\end{statement}

\emph{Proof of the Theorem from Proposition 2.} Let \(J\) be the Jacobian of a genus-\(2\) curve over
      \(\Q\), geometrically split, with a rational point of order \(13\), and suppose that \(J\) is simple over
      \(\Q\). Apply Proposition 2 to \(A=J\). If \(\Gamma\cong C_3\), the principal polarization of \(J\)
      contradicts (a). If \(\Gamma\cong C_2\), part (b) says that \(J\) with its canonical polarization is
      geometrically a product with the product polarization, which contradicts (F8). \qed

\emph{Proof of the Corollary.} If \(J\) is split over \(\Q\), it contains an elliptic curve
      \(E\) over \(\Q\), and by \hyperref[lemma-sequence]{Lemma~1} a rational point of order \(13\) on \(J\) gives one
      on \(E\) or on \(J/E\), contrary to (F9). If \(J\) is simple over \(\Q\), apply the Theorem. \qed

The proof of Proposition 2 occupies Section~\ref{type-n} to Section~\ref{c2}: type (N) in Section~\ref{type-n}, type
      (CM) in Section~\ref{type-cm}, \(\Gamma\in\{C_1,V_4,D_6,D_8,D_{12}\}\) in Section~\ref{noncyclic}, \(\Gamma=C_6\) in
      Section~\ref{c6}, \(\Gamma=C_4\) in Section~\ref{c4}, \(\Gamma=C_3\) in Section~\ref{c3}, and \(\Gamma=C_2\) in Section~\ref{c2}.
      In all cases except \(C_3\) and \(C_2\) we show that \(A(\Q)[13]=0\); no polarization hypothesis is
      needed there.

\section{Two non-isogenous geometric factors}\label{type-n}

Suppose that \(A\) is of type (N) and that \(P\in A(\Q)\) has order \(13\).

\begin{enumerate}
\item \(A_{\bar\Q}\) has exactly two elliptic subvarieties: for each \(E_i\), the space
          \(\Hom^0(E_i,A_{\bar\Q})\) has rank one over \(\End^0(E_i)\), so all nonzero maps
          \(E_i\to A_{\bar\Q}\) have the same image.
\item \(G_{\Q}\) permutes these two subvarieties. If it fixed them, they would be defined over \(\Q\)
          and \(A\) would not be simple over \(\Q\). Hence they are exchanged, and they are \(E\) and
          \(E^\sigma\) for an elliptic curve \(E\) over a quadratic field \(K\), where \(\sigma\) generates
          \(\Gal(K/\Q)\).
\item The quotient map restricts to a \(K\)-isogeny \(E^\sigma\to A_K/E\), so \(A_K/E\sim_K E^\sigma\).
\item By \hyperref[lemma-sequence]{Lemma~1}, either \(E(K)\) or \((A_K/E)(K)\) has a point of order \(13\). So there is
          an elliptic curve \(E''/K\) with a \(K\)-point of order \(13\) and \(E''\sim_K E\) or
          \(E''\sim_K E^\sigma\).
\item By Bosman--Bruin--Dujella--Najman~\cite[Theorem 11(1)]{BBDN}, an elliptic curve over a quadratic field with a point of order
          \(13\) is isomorphic over that field to its Galois conjugate: \(E''\cong_K E''^\sigma\).
\item Conjugating the isogeny of step 4 by \(\sigma\), we obtain \(E\sim_K E^\sigma\) in both cases, for
          instance \(E^\sigma\sim_K E''\cong_K E''^\sigma\sim_K E\) in the first case. This contradicts the
          assumption that \(E_1\) and \(E_2\) are not isogenous. \qed
\end{enumerate}

\section{Factors with complex multiplication}\label{type-cm}

Suppose that \(A\) is of type (CM), with CM field \(M\), and that \(A(\Q)[13]\neq0\).

\begin{statement}{Lemma 3.}\label{lemma-cm-field}
The field \(M\) is \(\Q(i)\) or \(\Q(\sqrt{-3})\).
\end{statement}

\emph{Proof.}

\begin{enumerate}
\item \emph{Good reduction over a totally ramified extension of \(\Q_2\).} The curve \(E\) has
          potentially good reduction everywhere, because it has CM (its \(j\)-invariant is an algebraic
          integer), and hence so does \(A\): good reduction
          over a finite extension passes to products and isogenies (Serre--Tate~\cite[§1, Theorem~1 and Corollaries~2 and~3]{ST}).
          By \cite[§2, Theorem~2(i)]{ST}, the inertia group \(I\subset G_{\Q_2}\) acts on \(T_3A\) through a finite
          quotient. Put \(N=I\cap\ker(\rho_{A,3})\), an open normal subgroup of \(I\) that is normal in
          \(G_{\Q_2}\). The group \(G_{\Q_2}/N\) is an extension of \(G_{\Q_2}/I\cong\hat\Z\) by the finite
          group \(I/N\); since \(\hat\Z\) is free, the extension splits, and there is a closed subgroup
          \(H\supseteq N\) with \(H\cap I=N\) and \(HI=G_{\Q_2}\). Its fixed field \(F\) is a finite extension
          of \(\Q_2\) with residue field \(\F_2\), i.e. totally ramified, whose inertia group \(N\) acts
          trivially on \(T_3A\). By the Néron--Ogg--Shafarevich criterion \cite[§1, Theorem~1]{ST}, \(A\) has good reduction
          \(\bar A/\F_2\) over \(F\). Prime-to-\(2\) torsion of \(A(F)\) injects into \(\bar A(\F_2)\), so
          \(13\) divides \(\#\bar A(\F_2)=P(1)\), where \(P\) is the characteristic polynomial of the
          Frobenius endomorphism \(\pi\) of \(\bar A\).
\item \emph{The Weil polynomial.} Write

\[
P(T)=(T^2-\beta_1T+2)(T^2-\beta_2T+2)=T^4+aT^3+bT^2+2aT+4
\]

with real \(|\beta_i|\le2\sqrt2\). The script \href{https://genus-2-torsion.github.io/data/knowledge/sources/qsimple_13_cm_weil.m}{\texttt{qsimple\_\allowbreak{}13\_\allowbreak{}cm\_\allowbreak{}weil.m}} enumerates all such
          integer pairs \((a,b)\) exactly; the values of \(P(1)\) are \(1,\dots,16,19,20,25\), and the only
          polynomial with \(13\mid P(1)\) is

\[
P(T)=T^4+2T^3+2T^2+4T+4,\qquad P(1)=13.
\]

It is irreducible, and its root field is \(\Q(\zeta_{12})\).
\item \emph{Frobenius commutes with \(M\).} All geometric endomorphisms of \(A\) are defined over a
          finite extension \(F'\) of \(F\) and reduce injectively into \(\End^0(\bar A_{\bar\F_2})\), compatibly
          with the Galois action: if \(g\in G_F\) induces the Frobenius of the residue field, then the
          reduction of \(g(m)\) is \(\pi\bar m\pi^{-1}\). The field \(M\) is the centre of
          \(\End^0(A_{\bar\Q})=M_2(M)\), so \(g(m)\in\{m,\bar m\}\) for \(m\in M\). If \(g\) acted on \(M\) by
          complex conjugation, then \(\pi\bar s\pi^{-1}=-\bar s\) for the reduction \(\bar s\) of
          \(s=\sqrt{-d}\in M\), so \(\pi\) and \(-\pi\) would be conjugate and \(P(T)=P(-T)\), i.e. \(a=0\).
          But \(a=2\). Hence \(\pi\) commutes with the reduction of \(M\).
\item \emph{Conclusion.} The endomorphisms of \(\bar A_{\bar\F_2}\) commuting with \(\pi\) are those defined
          over \(\F_2\). Since \(P\) has no multiple root, \(\End^0_{\F_2}(\bar A)=\Q(\pi)\cong\Q(\zeta_{12})\)
          by Tate~\cite[Theorem 2(c)]{Tat}. So \(M\) embeds into \(\Q(\zeta_{12})\), whose imaginary quadratic
          subfields are \(\Q(i)\) and \(\Q(\sqrt{-3})\). \qed
\end{enumerate}

Let \(E_M/\Q\) be \(y^2=x^3-x\) if \(M=\Q(i)\), and \(y^2=x^3+1\) if \(M=\Q(\sqrt{-3})\). It has CM by
      \(\mathcal O_M\), defined over \(M\), and good ordinary reduction at \(13\), which splits in \(M\)
      (conductors \(32\) and \(36\); \href{https://genus-2-torsion.github.io/data/knowledge/sources/qsimple_13_cm_weil.m}{\texttt{qsimple\_\allowbreak{}13\_\allowbreak{}cm\_\allowbreak{}weil.m}}).

\begin{statement}{Proposition 4.}\label{prop-cm}
If \(A\) is of type (CM), then \(A(\Q)[13]=0\).
\end{statement}

\emph{Proof.} Suppose \(A(\Q)[13]\neq0\); by \hyperref[lemma-cm-field]{Lemma~3}, \(M\) is \(\Q(i)\) or
      \(\Q(\sqrt{-3})\).

\begin{enumerate}
\item \emph{The twisting representation.} Since \(A_{\bar\Q}\sim E_M^2\), the surface \(A\) is a form of
          \(E_M^2\) in the isogeny category, given by a cocycle \(g\mapsto a_g\in\GL_2(M)=\operatorname{Aut}^0(E^2_{M,\bar\Q})\),
          where \(G_{\Q}\) acts on \(\GL_2(M)\) through \(\Gal(M/\Q)\). On \(G_M\) the cocycle is a continuous
          homomorphism \(\rho:G_M\to\GL_2(M)\), so it has finite image. An element of \(\GL_2(M)\) of finite
          order \(n\) has eigenvalues of degree at most \(2\) over \(M\), so \(\varphi(n)\le4\) and
          \(13\nmid n\); hence \(13\nmid\#\rho(G_M)\). Transport by a quasi-isogeny shows that
          \(g\in G_M\) acts on \(V_\ell A\cong V_\ell E_M\otimes_{M\otimes\Q_\ell}(M\otimes\Q_\ell)^2\) by
          \(\rho_{E_M,\ell}(g)\otimes\rho(g)\).
\item \emph{The two \(\mathfrak p\)-components.} Write \(13=\mathfrak p\bar{\mathfrak p}\) in \(M\), with
          embeddings \(\iota_{\mathfrak p},\iota_{\bar{\mathfrak p}}:M\to\Q_{13}\), and
          \(V_{13}E_M=V_{\mathfrak p}\oplus V_{\bar{\mathfrak p}}\), on which \(G_M\) acts by characters
          \(\alpha_{\mathfrak p},\alpha_{\bar{\mathfrak p}}:G_M\to\Z_{13}^\times\). Then
          \(V_{13}A=V_{\mathfrak p}(A)\oplus V_{\bar{\mathfrak p}}(A)\), with
          \(V_{\mathfrak p}(A)\cong\alpha_{\mathfrak p}\otimes\iota_{\mathfrak p}\rho\) as \(G_M\)-modules. This is the
          decomposition under the centre \(M\) of \(\End^0(A_{\bar\Q})\); since inner automorphisms act
          trivially on the centre, \(G_{\Q}\) acts on \(M\) through \(\Gal(M/\Q)\).
\item \emph{Duality.} Fix a polarization of \(A\) over \(\Q\). Its Rosati involution is positive, so it
          is complex conjugation on the centre \(M\), and the Riemann form satisfies
          \(E(mx,y)=E(x,\bar my)\). For \(m\in M\setminus\Q\) and \(x,y\in V_{\mathfrak p}(A)\) this gives
          \((\iota_{\mathfrak p}(m)-\iota_{\mathfrak p}(\bar m))E(x,y)=0\), so \(V_{\mathfrak p}(A)\) is isotropic, and
          \(E\) pairs it perfectly with \(V_{\bar{\mathfrak p}}(A)\): \(V_{\bar{\mathfrak p}}(A)\cong V_{\mathfrak p}(A)^\vee(1)\)
          as \(G_M\)-modules.
\item \emph{Complex conjugation.} Complex conjugation \(c\) satisfies \(cm=\bar mc\) on \(V_{13}A\), so it
          maps \(V_{\mathfrak p}(A)\) onto \(V_{\bar{\mathfrak p}}(A)\) and \(V_{\bar{\mathfrak p}}(A)\cong V_{\mathfrak p}(A)^c\).
\item \emph{Both residual pieces are \(\mathbf 1\oplus\bar\chi_{13}\).} A rational point of order \(13\) gives a
          trivial \(G_M\)-constituent of \(A[13]\), hence, by (F4) applied to the lattice
          \((T_{13}A\cap V_{\mathfrak p}(A))\oplus(T_{13}A\cap V_{\bar{\mathfrak p}}(A))\), of
          \(\bar V_{\mathfrak p}^{\mathrm{ss}}\) or \(\bar V_{\bar{\mathfrak p}}^{\mathrm{ss}}\). By step 4 both contain
          \(\mathbf 1\). Write \(\bar V_{\mathfrak p}^{\mathrm{ss}}=\mathbf 1\oplus\beta\); by step 3,
          \(\bar V_{\bar{\mathfrak p}}^{\mathrm{ss}}=\bar\chi_{13}\oplus\beta^{-1}\bar\chi_{13}\), and since
          \(\bar\chi_{13}|_{G_M}\neq1\) we get \(\beta=\bar\chi_{13}\):

\[
\bar V_{\mathfrak p}(A)^{\mathrm{ss}}\cong\bar V_{\bar{\mathfrak p}}(A)^{\mathrm{ss}}\cong\mathbf 1\oplus\bar\chi_{13}\qquad\text{as }G_M\text{-modules}.
\]

Therefore \(\bar\rho^{\mathrm{ss}}\cong\bar\alpha_{\mathfrak p}^{-1}\oplus\bar\alpha_{\mathfrak p}^{-1}\bar\chi_{13}\),
          where \(\bar\rho\) is the reduction of \(\iota_{\mathfrak p}\rho\), which is defined because \(\rho\) has
          finite image of order prime to \(13\).
\item \emph{Inertia above \(13\).} At each place \(v\mid13\) of \(M\) we have \(M_v=\Q_{13}\), and \(E_M\) has
          good ordinary reduction. The connected part of \(E_M[13^\infty]\) at \(v\) is \(\mathcal O_M\)-stable,
          hence equal to \(E_M[\mathfrak p^\infty]\) or \(E_M[\bar{\mathfrak p}^\infty]\); the inertia group \(I_v\)
          acts on it by \(\chi_{13}\) and trivially on the étale quotient. Since
          \(\alpha_{\bar{\mathfrak p}}(g)=\alpha_{\mathfrak p}(c^{-1}gc)\) and \(c\) exchanges the two places
          above \(13\), there is a place \(w\mid13\) with \(\alpha_{\mathfrak p}|_{I_w}=\chi_{13}|_{I_w}\).
          There, \(\bar\rho^{\mathrm{ss}}|_{I_w}\cong\bar\chi_{13}^{-1}\oplus\mathbf 1\).
\item \emph{Contradiction.} Choose \(g_0\in I_w\) such that \(\bar\chi_{13}(g_0)\) generates
          \(\F_{13}^\times\). The eigenvalues of \(\iota_{\mathfrak p}\rho(g_0)\) are roots of unity of order prime
          to \(13\), on which reduction is injective, so they are the Teichmüller lifts of
          \(\bar\chi_{13}(g_0)^{-1}\) and \(1\): a primitive \(12\)-th root of unity and \(1\). Hence
          \(\det\rho(g_0)\in M\) is a primitive \(12\)-th root of unity. But \(\Q(i)\) and \(\Q(\sqrt{-3})\)
          contain only \(4\) and \(6\) roots of unity. \qed
\end{enumerate}

\section{Trivial or non-cyclic endomorphism group}\label{noncyclic}

In this section \(A\) is of type (S), \(p\) is an odd prime, \(V=V_p(A)\) and \(T=T_p(A)\). The group
      \(G_{\Q}\) acts on \(D\otimes\Q_p\) through \(\Gamma\), and \(g(xt)=g(x)\,g(t)\) for \(x\in D\otimes\Q_p\) and
      \(t\in V\).

\begin{statement}{Lemma 5 (polarization).}\label{lemma-a}
Let \(T'\subset V\) be a \(G_{\Q}\)-stable lattice such that \((T'/pT')|_{G_L}\) has trivial
          semisimplification. Then \(\Q(\zeta_p)\subseteq L\).
\end{statement}

\emph{Proof.} A polarization of \(A\) over \(\Q\) gives \(V\cong V^\vee\otimes\Q_p(1)\). Transport the
      lattice \(\Hom(T',\Z_p)\otimes\Z_p(1)\) into \(V\). Its reduction has semisimplification
      \((\bar\chi_p|_{G_L})^{4}\) as a \(G_L\)-module, while \(T'/pT'\) has \(\mathbf 1^4\). By (F4) over
      \(G_L\), \(\bar\chi_p|_{G_L}=1\). The degree of the polarization plays no role. \qed

\begin{statement}{Lemma 6 (unramified order).}\label{lemma-b}
Suppose that \(\Gamma\) normalizes an unramified quadratic \(\Z_p\)-order \(B\subset D\otimes\Q_p\)
          (the ring of integers of an unramified quadratic extension of \(\Q_p\)), that some \(g\in\Gamma\) is
          represented by a unit \(R_g\in B^\times\) whose reduction lies in \(\F_{p^2}\setminus\F_p\), and that
          \(A(\Q)[p]\neq0\). Then \(\Q(\zeta_p)\subseteq L\).
\end{statement}

\emph{Proof.}

\begin{enumerate}
\item \emph{A \(\Gamma\)-stable maximal order.} The \(B\)-stable lattices in \(\Q_p^2\) are the fractional
          ideals of \(B\otimes\Q\) times a fixed vector, so they form one homothety class \([\Lambda]\). For
          \(g\in\Gamma\), \(R_g\Lambda\) is again \(B\)-stable, so \(\Gamma\) preserves \([\Lambda]\) and the maximal
          order \(\mathcal R=\End(\Lambda)\cong M_2(\Z_p)\).
\item \emph{Tensor decomposition.} Put \(T'=\mathcal R T\), a \(G_{\Q}\)-stable lattice. By the argument of (F6), with \(e_{11}T'\) now of rank \(2\) over \(\Z_p\), \(T'=e_{11}T'\oplus e_{21}e_{11}T'\cong\Lambda\otimes W_0\) with \(W_0=e_{11}T'\). Let \(\bar V=T'/pT'\) and \(U=\Lambda/p\Lambda\). The action of \(G_{\Q}\)
          on \(\bar V\) normalizes the action of \(M_2(\F_p)=\End(U)\), and its projective action on \(U\)
          factors through \(\Gamma\). Pulling back \(1\to\F_p^\times\to\GL(U)\to\PGL(U)\to1\) along it gives a
          group \(\tilde G\) and two-dimensional \(\tilde G\)-modules \(U,W\) with \(\bar V\cong U\otimes W\); an
          element \((g,M)\) acts on \(U\) by \(M\), and \((M^{-1}\otimes1)\circ g\) commutes with \(\End(U)\), so
          it acts on the second factor.
\item \emph{\(U\) is irreducible.} For the element \(g\) of the hypothesis, \(R_g\in B^\times\subset\GL(\Lambda)\)
          induces the action of \(g\) on \(\mathcal R\), so \(\tilde G\) contains elements acting on \(U\) by the
          reduction of \(R_g\). That is multiplication by an element of \(\F_{p^2}\setminus\F_p\) on the
          one-dimensional \(\F_{p^2}\)-space \(U\), which has no eigenvector in \(U\).
\item \emph{\(W\) is irreducible.} By (F4) and \(A(\Q)[p]\neq0\), \(\bar V\) has a trivial constituent. If
          \(W\) had a stable line \(\ell\), then \(U\otimes\ell\) and \(U\otimes W/\ell\) would be irreducible of
          dimension \(2\), and \(\bar V\) would have no trivial constituent.
\item \emph{\(\bar V|_{G_L}\) is trivial.} The subgroup \(H=\{(h,1):h\in G_L\}\) is normal in \(\tilde G\)
          and acts trivially on \(U\). By Clifford's theorem, \(W|_H\) is semisimple. Since
          \(\bar V|_H\cong W|_H^{\oplus2}\) has a trivial constituent, \(W^H\neq0\); as \(W^H\) is
          \(\tilde G\)-stable and \(W\) is irreducible, \(W^H=W\).
\item Lemma 5 now gives \(\Q(\zeta_p)\subseteq L\). \qed
\end{enumerate}

\begin{statement}{Lemma 7 (split blocks).}\label{lemma-c}
Suppose that \(\Gamma\) normalizes a split étale order \(B=\Z_pe_1\oplus\Z_pe_2\subset D\otimes\Q_p\),
          and that some element of \(\Gamma\) exchanges \(e_1\) and \(e_2\). Let \(K\subset L\) be the fixed field
          of the index-\(2\) subgroup fixing \(e_1\). Suppose that \(G_K\) acts on the line
          \(e_1(D\otimes\Q_p)e_2\) by a character \(\eta\) that factors through \(\Gal(L/K)\) and has
          \(\bar\eta\neq1\). If \(A(\Q)[p]\neq0\), then \(\Q(\zeta_p)\subseteq L\).
\end{statement}

\emph{Proof.} Put \(T'=BT=T_1\oplus T_2\) with \(T_i=e_iT'\), and \(M_i=T_i/pT_i\), \(V_i=T_i\otimes\Q_p\).
      The group \(G_K\) preserves each \(T_i\), and elements of \(G_{\Q}\setminus G_K\) exchange them. A
      nonzero \(X\in e_1(D\otimes\Q_p)e_2\) maps \(V_2\) isomorphically onto \(V_1\) and satisfies
      \(g(Xv)=\eta(g)Xg(v)\) for \(g\in G_K\), so \(V_1\cong\eta\otimes V_2\) and, by (F4),
      \(M_1^{\mathrm{ss}}\cong\bar\eta\otimes M_2^{\mathrm{ss}}\). The rational \(p\)-torsion point gives a trivial
      \(G_K\)-constituent of \(M_1\oplus M_2\); conjugation by an element outside \(G_K\) exchanges the two
      blocks and preserves the trivial character, so both \(M_i\) have a trivial constituent. Writing
      \(M_2^{\mathrm{ss}}=\mathbf 1\oplus\alpha\), we get \(M_1^{\mathrm{ss}}=\bar\eta\oplus\bar\eta\alpha\ni\mathbf 1\),
      so \(\alpha=\bar\eta^{-1}\). All four constituents are trivial on \(G_L\), and Lemma 5 applies. \qed

We apply these lemmas with \(p=13\), assuming \(A(\Q)[13]\neq0\). Note that \(-1\equiv5^2\),
      \(-3\equiv6^2\) and \(-1/3\equiv2^2\pmod{13}\).

\begin{itemize}
\item \textbf{\(\Gamma=C_1\).} Then \(\End^0_{\Q}(A)=D=M_2(\Q)\), and \(A\) is not simple over \(\Q\).
\item \textbf{\(\Gamma=D_{2n}\), \(n\in\{3,4,6\}\).} Let \(R\) be the trace-one representative of a rotation
          \(r\) of order \(n\). A reflection conjugates \(r\) to \(r^{-1}\), whose trace-one representative is
          \(1-R\) (as \(R(1-R)=1/c\)); so \(B=\Z_{13}[R]\) is \(\Gamma\)-stable. Its discriminant
          \(-3,-1,-1/3\) is a nonzero square modulo \(13\), so \(B\) is split, with idempotents fixed by the
          rotations and exchanged by the reflections. For \(X\in e_1De_2\), \(\mathrm{Ad}(R)X=(\lambda_1/\lambda_2)X\),
          where \(\lambda_1/\lambda_2\) is a primitive \(n\)-th root of unity in \(\Z_{13}^\times\); its reduction
          still has order \(n\). Lemma 7 gives \(\Q(\zeta_{13})\subseteq L\), so \(\Gamma\) surjects onto
          \(\Gal(\Q(\zeta_{13})/\Q)\cong C_{12}\). But the abelianization of \(D_{2n}\) has exponent \(2\).
\item \textbf{\(\Gamma=V_4\).} The map \([M]\mapsto v_{13}(\det M)\bmod2\) is a well-defined homomorphism
          \(\Gamma\to\Z/2\); choose \(r\neq1\) in its kernel. A representative \(R\) of \(r\) has trace \(0\)
          and \(R^2=a\in\Q^\times\) with \(v_{13}(a)\) even; scale \(R\) so that \(a\) is a \(13\)-adic unit.
          Since \(\Gamma\) is abelian, every \(g\in\Gamma\) sends \(R\) to \(\pm R\), so \(B=\Z_{13}[R]\) is a
          \(\Gamma\)-stable étale order. If \(a\) is not a square modulo \(13\), \(B\) is unramified and \(R\) is a
          unit of \(B\) with reduction in \(\F_{169}\setminus\F_{13}\); Lemma 6 applies. If \(a\) is a square,
          \(B\) is split. Some \(g\in\Gamma\) sends \(R\) to \(-R\), since otherwise \(\Gamma\) would lie in the
          torus \((\Q_{13}[R])^\times/\Q_{13}^\times\cong\Q_{13}^\times\), whose finite subgroups are cyclic; this
          \(g\) exchanges the idempotents, the subgroup fixing them is \(\{1,r\}\), and \(r\) acts on
          \(e_1De_2\) by \(-1\). Lemma 7 applies. In both cases \(\Q(\zeta_{13})\subseteq L\), which is
          impossible since \([L:\Q]=4\).
\end{itemize}

Hence, if \(A(\Q)[13]\neq0\), then \(\Gamma\in\{C_2,C_3,C_4,C_6\}\).

\section{The case \(\Gamma=C_6\)}\label{c6}

In this section \(A_{\bar\Q}\sim E^2\) with \(E\) without CM and \(\Gamma=\langle r\rangle\cong C_6\). Let
      \(R\) be the trace-one representative of \(r\), so \(R^2-R+\tfrac13=0\), and let \(K=L^{\langle r^2\rangle}\)
      be the quadratic subfield of \(L\).

\textbf{\(L\) is totally real.} Fix an embedding \(\bar\Q\subset\mathbb C\), let \(c\) be complex
      conjugation and \(H=H_1(A(\mathbb C),\mathbb R)\). As a module over \(D_{\mathbb R}=M_2(\mathbb R)\),
      \(H=U\otimes W\) with \(U=\mathbb R^2\) the standard module. The complex structure commutes with \(D\),
      so it is \(1\otimes J\) with \(J^2=-1\). The map \(c_*\) satisfies \(c_*\phi=c(\phi)c_*\) for \(\phi\in D\),
      where \(c\) acts on \(D\) by \(\mathrm{Ad}(M)\) for some \(M\in\GL_2(\mathbb R)\); hence
      \(c_*=M\otimes N\). Antiholomorphy gives \(NJ=-JN\), so in a suitable basis
      \(N=\left(\begin{smallmatrix}a&b\\b&-a\end{smallmatrix}\right)\) and \(N^2=(a^2+b^2)I\). From
      \(c_*^2=1\), \(M^2\) is a positive scalar, so either \(M\) is scalar, or \(M\) has eigenvalues
      \(\pm\mu\) with \(\mu\) real and \(\det M<0\). Every element of \(\Gamma\) is \([R^k]\) with
      \(\det(\lambda R^k)=\lambda^2 3^{-k}>0\) for real \(\lambda\). Hence \(c\) acts trivially on \(D\), i.e.
      \(c\in G_L\). As this holds for every embedding, \(L\) is totally real; in particular \(K\) is real.

\textbf{Normalization.} \(\End^0_{\Q}(A)=D^\Gamma=\Q[R]\cong\Q(\sqrt{-3})\). The element \(u=3R-1\)
      satisfies \(u^2-u+1=0\), and \(R=(1+u)/3\), \(R^2=u/3\). Put \(\mathcal O=\Z[u]\) and choose \(m\ge1\)
      with \(mu\in\End_{\Q}(A)\). The image \(A'\) of \(a\mapsto(ma,mua)\) in \(A^2\) is \(\Q\)-isogenous to
      \(A\), and \(u\) acts on it by the integral endomorphism \((x,y)\mapsto(y,y-x)\) of \(A^2\). So
      \(\mathcal O\subset\End_{\Q}(A')\) and \(A'[13]^{\mathrm{ss}}\cong A[13]^{\mathrm{ss}}\).

\textbf{A twist.} Write \(g|_L=r^{k(g)}\) and \(\xi(g)=u^{-2k(g)}\in\mu_3\subset\operatorname{Aut}_{\Q}(A')\);
      this is a homomorphism, hence a cocycle. Let \(B/\Q\) be the twist, with
      \(f:A'_{\bar\Q}\to B_{\bar\Q}\) and \(f^{-1}\circ g(f)=\xi(g)\). Transported by \(f\), the action of \(g\)
      on \(D\) becomes \(\mathrm{Ad}(\xi(g)R^{k(g)})=\mathrm{Ad}(S^{k(g)})\) with \(S=u^{-2}R\), and
      \(S^2=u^{-4}\cdot u/3\) \({}=-1/3\) is scalar. Put \(T=3S=u^{-2}(1+u)\in\mathcal O\); then \(T^2=-3\). Since
      \(\mathcal O\) commutes with \(\xi\), we have \(\mathcal O\subset\End_{\Q}(B)\); all geometric
      endomorphisms of \(B\) are defined over \(K\), and \(\tau(\phi)=T\phi T^{-1}\) for
      \(\tau\in G_{\Q}\setminus G_K\). On Tate modules, \(f\) identifies the action of \(g\) on \(T_\ell B\) with
      \(x\mapsto\xi(g)\,g(x)\) on \(T_\ell A'\).

\textbf{A \(3\)-isogeny.} Choose a rank-one idempotent \(e\in D\) and \(n\) with
      \(ne\in\End(B_K)\), and put \(E_0=ne(B)\), an elliptic curve over \(K\). For \(\tau\notin G_K\),
      \(\tau(E_0)\) is the image of \(T(ne)T^{-1}\), so \(\tau(E_0)=T(E_0)\). Moreover \(T(E_0)\neq E_0\),
      since otherwise \(T|_{E_0}\) would be an endomorphism with square \(-3\). The isogenies
      \(T|_{E_0}:E_0\to\tau E_0\) and \(T|_{\tau E_0}:\tau E_0\to E_0\) are Galois conjugate, hence have the same
      degree \(d\), and their composite is \([-3]\); so \(d=3\). Thus \(T|_{\tau E_0}:\tau(E_0)\to E_0\) is a
      cyclic \(K\)-isogeny of degree \(3\).

\begin{statement}{Proposition 8.}\label{prop-c6}
Let \(A/\Q\) be an abelian surface with \(A_{\bar\Q}\sim E^2\), \(E\) without CM, and
          \(\Gamma\cong C_6\). Then \(A[13]^{\mathrm{ss}}\) has no one-dimensional constituent. In particular
          \(A(\Q)[13]=0\).
\end{statement}

\emph{Proof.} Suppose that \(A[13]^{\mathrm{ss}}\) has a one-dimensional constituent.

\begin{enumerate}
\item \emph{The \(\mathfrak l\)-block.} Since \(-3\equiv6^2\pmod{13}\),
          \(\mathcal O\otimes\Z_{13}=\Z_{13}\times\Z_{13}\), with primes \(\mathfrak l,\bar{\mathfrak l}\), and
          \(T_{13}A'=T_{\mathfrak l}\oplus T_{\bar{\mathfrak l}}\) with both summands \(G_{\Q}\)-stable of rank \(2\).
          Relabel so that \(T_{\mathfrak l}(A')/13\) has a one-dimensional constituent.
\item \emph{Transfer to \(B\).} By the twist, \(T_{\mathfrak l}(B)\cong T_{\mathfrak l}(A')\otimes\xi_{\mathfrak l}\)
          with \(\xi_{\mathfrak l}:G_{\Q}\to\mu_3(\Z_{13})\). So \(T_{\mathfrak l}(B)/13\) is reducible and has a
          \(G_{\Q}\)-stable line \(\bar C\).
\item \emph{\(V_{13}E_0\) is the \(\mathfrak l\)-block.} The map
          \(\iota:V_{13}E_0\to V_{13}B\to V_{\mathfrak l}(B)\) is \(G_K\)-equivariant. By (F10), \(V_{13}E_0\) is
          irreducible, so \(\iota\) is injective or zero. If it were zero, \(V_{13}E_0=V_{\bar{\mathfrak l}}(B)\)
          would be \(T\)-stable, so the abelian subvariety \(E_0+T(E_0)\) would have the same \(V_{13}\) as
          \(E_0\), forcing \(T(E_0)=E_0\). Hence \(\iota\) is an isomorphism.
\item \emph{The curve \(P\).} By (F1), the lattice \(\iota^{-1}(T_{\mathfrak l}B)\) gives an elliptic curve
          \(P/K\), a \(K\)-isogeny \(h:E_0\to P\) of \(13\)-power degree, and a \(G_K\)-isomorphism
          \(\iota_P:T_{13}P\to T_{\mathfrak l}B\). Put \(C=\iota_P^{-1}(\bar C)\), a \(G_K\)-stable subgroup of
          order \(13\) of \(P\).
\item \emph{An isogeny \(P^\tau\to P\).} Fix \(\tau\notin G_K\) and let
          \(\Phi=\iota_P^{-1}\circ\tau_B\circ\iota_P\circ\tau_*^{-1}\) \({}:T_{13}(P^\tau)\to T_{13}P\), where \(\tau_B\) is
          the action of \(\tau\) on \(T_{\mathfrak l}B\). Since \(\tau^{-1}g\tau\in G_K\) for \(g\in G_K\) and
          \(\iota_P\) is \(G_K\)-linear, \(\Phi\) is \(G_K\)-equivariant. By (F10),
          \(\Phi=\lambda\nu_0\), where \(\nu_0\) generates \(\Hom_K(P^\tau,P)\cong\Z\) (nonzero via \(T\)) and
          \(\lambda\in\Z_{13}\). As \(\Phi\) is invertible, \(\lambda\in\Z_{13}^\times\) and
          \(13\nmid\deg\nu_0\). A generator is a cyclic isogeny.
\item \emph{\(\deg\nu_0=3\).} The \(K\)-isogeny \(h\circ T|_{\tau E_0}\circ\widehat{h^\tau}:P^\tau\to P\) has
          degree \(3\cdot13^{2a}\) and equals \(n\nu_0\) for some \(n\in\Z\). Since \(13\nmid\deg\nu_0\) and
          \(3\cdot13^{2a}\) is not a square, \(\deg\nu_0=3\).
\item \emph{\(\nu_0\) respects the \(13\)-structure.} Since \(\tau_B(\bar C)=\bar C\), we get
          \(\Phi(\tau C)=C\), hence \(\nu_0(\tau C)=C\).
\item \emph{A point of \(Y_0(39)(K)\).} Put \(\nu_0^\tau=\tau(\nu_0):P\to P^\tau\). Then
          \(\nu_0\nu_0^\tau\in\End_K(P)=\Z\) has degree \(9\), so it is \(\pm3\) and
          \(\nu_0^\tau(P[3])=\ker\nu_0\). Applying \(\tau\) to \(\nu_0(\tau C)=C\) gives \(\nu_0^\tau(C)=\tau C\).
          The point \(x=(P,\ \ker\nu_0^\tau+C)\in Y_0(39)(K)\) satisfies

\[
\tau(x)=(P^\tau,\ \ker\nu_0+\tau C)=w_3(x),
\]

using the isomorphism \(P/\ker\nu_0^\tau\cong P^\tau\) induced by \(\nu_0^\tau\). The isogeny
          \(\nu_0^\tau\) of degree \(3\) carries the \(13\)-structure \(C\) to \(\tau C\); this is the moduli
          action of \(w_3\).
\item \emph{\(x\) is not rational.} If \(\tau(x)=x\), then \(w_3(x)=x\), and \(P\) would have a cyclic
          endomorphism of degree \(3\), hence CM; but \(P\sim E_0\).
\item \emph{Contradiction.} The image of \(x\) in \(X_0(39)/w_3\) is fixed by \(\Gal(K/\Q)\), hence rational.
          By \href{https://genus-2-torsion.github.io/data/knowledge/sources/qsimple_13_x0_39.m}{\texttt{qsimple\_\allowbreak{}13\_\allowbreak{}x0\_\allowbreak{}39.m}}, \(X_0(39)/w_3\) is the elliptic curve 39a1, of rank \(0\) with torsion
          \(\Z/2\times\Z/2\); above two of its four rational points lie the four rational cusps of
          \(X_0(39)\), and above each of the other two lies a pair of conjugate points with residue field
          \(\Q(\sqrt{-7})\). The non-rational point \(x\) therefore has residue field \(\Q(\sqrt{-7})\subseteq K\),
          so \(K=\Q(\sqrt{-7})\), contradicting the fact that \(K\) is real. \qed
\end{enumerate}

\section{The case \(\Gamma=C_4\)}\label{c4}

In this section \(A\) is of type (S), \(\Gamma=\langle r\rangle\cong C_4\), and \(K\) is the
      quadratic subfield of \(L\). The trace-one representative \(R\) of \(r\) has \(\det R=\tfrac12\). Put
      \(J=2R-1\); then \(J^2=-1\), \(R=(1+J)/2\) and \(R^2=J/2\). So \(\Gal(L/K)=\langle r^2\rangle\) acts on \(D\) by
      \(\mathrm{Ad}(J)\), and \(\End^0_K(A_K)\) is the centralizer \(\Q(J)\cong\Q(i)\) of \(J\), a field: \textbf{\(A_K\)
      is simple over \(K\)}. We show that \(A(\Q)[13]=0\).

\begin{statement}{Lemma 9 (degree class).}\label{lemma-degree}
Let \(X\) be an elliptic curve without CM over a number field \(L_1\supseteq L\), and
          \(h:X^2\to A_{L_1}\) a quasi-isogeny defined over \(L_1\). For \(g\in G_{\Q}\) let \(\mu_g\) be a
          generator of \(\Hom(X^g_{\bar\Q},X_{\bar\Q})\cong\Z\). If \(\det R_g>0\), then
          \(\deg\mu_g\equiv\det R_g\pmod{\Q^{\times2}}\).
\end{statement}

\emph{Proof.} Put \(N_g=h^{-1}\circ g(h)\) \(\in\Hom^0((X^g)^2,X^2)=M_2(\Q)\mu_g\) and write \(N_g=\mu_gN_0\)
      with \(N_0\in\GL_2(\Q)\). For \(Y\in M_2(\Z)=\End(X^2)\), \(g(hYh^{-1})=g(h)Yg(h)^{-1}\) \({}=hN_0YN_0^{-1}h^{-1}\),
      so in the identification \(D\cong M_2(\Q)\) given by \(h\), \(g\) acts by \(\mathrm{Ad}(N_0)\). Hence
      \([N_0]\) is conjugate to \([R_g]\) in \(\PGL_2(\Q)\), and \(\det N_0\equiv\det R_g\) modulo
      \(\Q^{\times2}\); in particular \(\det N_0>0\). Since \(\deg g(h)=\deg h\), (F3) gives
      \(1=\deg N_g=(\deg\mu_g)^2(\det N_0)^2\), so \(\det N_0=1/\deg\mu_g\). \qed

\begin{statement}{Lemma 10.}\label{lemma-two-subgroup}
Let \(X\) be an elliptic curve without CM. Two distinct subgroups \(D_1,D_2\subset X[2]\) of order
          \(2\) have non-isomorphic quotients \(X/D_1\), \(X/D_2\).
\end{statement}

\emph{Proof.} Let \(\pi_i:X\to X/D_i\) and suppose \(\phi:X/D_1\to X/D_2\) is an isomorphism. Then
      \(\hat\pi_2\phi\pi_1\) is an endomorphism of degree \(4\), hence \(\pm2=\pm\hat\pi_2\pi_2\). So
      \(\hat\pi_2\circ(\phi\pi_1\mp\pi_2)=0\), and \(\phi\pi_1=\pm\pi_2\) because \(X\) is connected and
      \(\ker\hat\pi_2\) is finite. Hence \(D_1=D_2\). \qed

Suppose now that \(A(\Q)[13]\neq0\).

\textbf{Step 1: a curve over \(K\).}

\begin{enumerate}
\item Let \(E\subset A_L\) be an elliptic subvariety and \(\sigma\) the nontrivial element of \(\Gal(L/K)\).
          Then \(E^\sigma\neq E\), since otherwise \(E\) would be defined over \(K\). So
          \(A_L\sim_L E\times E^\sigma\sim_L E^2\).
\item Apply \hyperref[lemma-degree]{Lemma~9} with \(X=E\), \(L_1=L\), and a lift \(g\) of \(\sigma\): since
          \(\det R_g=\det R^2=\tfrac14\), a generator \(\mu:E^\sigma\to E\) of \(\Hom(E^\sigma,E)\) has square
          degree \(m^2\). It is defined over \(L\) (Section~\ref{cases}), and its kernel is cyclic, since \(\mu\) is
          primitive.
\item \emph{The middle curve.} Let \(H_1\subset\ker\mu\) be the subgroup of order \(m\), which is
          \(G_L\)-stable, and \(E_c=E^\sigma/H_1\). Then \(\mu\) factors as \(E^\sigma\to E_c\to E\) through two cyclic
          isogenies of degree \(m\). The conjugate \(\sigma(\mu):E\to E^\sigma\) and the dual \(\hat\mu\) both
          generate \(\Hom(E,E^\sigma)\cong\Z\), so \(\sigma(\mu)=\pm\hat\mu\). The unique subgroup of order \(m\)
          of \(\ker\sigma(\mu)=\ker\hat\mu\) is both \(\sigma(H_1)\) and the kernel of
          \(\hat\beta:E\to E_c\), where \(\beta:E_c\to E\) is the second factor. Hence
          \(E_c^\sigma=E/\sigma(H_1)\cong E_c\), and \(j(E_c)\in K\).
\item Let \(E_K\) be a model of \(E_c\) over \(K\); it exists, and \(j(E_K)\neq0,1728\). By (F7) over \(K\),
          \(A_K\sim_K E_K\otimes\rho_K\), and the projective image of \(\rho_K\) is \(\langle[J]\rangle\cong C_2\).
\item The finite subgroups of \(\GL_2(\Q)\) with projective image of order \(2\) are \(\{1,g\}\) and
          \(\{\pm1,\pm g\}\) with \(g^2=1\), \(g\neq\pm1\), and \(\langle g\rangle\) with \(g^2=-1\). In the first two
          cases \(g\) is diagonalizable over \(\Q\) and \(A_K\) would split over \(K\). So
          \(\rho_K(G_K)\cong C_4\) and \(A_K\sim_K E_K\otimes\Z[i](\psi_K)\) for a character
          \(\psi_K:G_K\to\mu_4\) of order \(4\).
\end{enumerate}

\textbf{Step 2: a \(K\)-rational \(2\)-isogeny to the conjugate.}

\begin{enumerate}
\item Choose \(\tau\in G_{\Q}\) with \(\tau|_K\neq\mathrm{id}\); then \(\tau|_L\in\{r,r^3\}\) and
          \(\tau^2\in G_K\), so \(E_K^{\tau^2}=E_K\). Over the kernel field \(L_1\) of \(\psi_K\) we have
          \(A_{L_1}\sim_{L_1}E_K^2\) with all endomorphisms defined over \(L_1\), so \(L_1\supseteq L\).
          \hyperref[lemma-degree]{Lemma~9} with \(\det R_\tau\in\{\tfrac12,\tfrac18\}\) shows that a generator
          \(\mu:E_K^\tau\to E_K\) of \(\Hom(E^\tau_{K,\bar\Q},E_{K,\bar\Q})\) has degree \(2n^2\).
\item Every \(g\in G_K\) maps \(\mu\) to \(\pm\mu\). So the cyclic group \(\ker\mu\) and all its subgroups
          are \(K\)-rational, and \(\mu\) factors as
          \(E_K^\tau\xrightarrow{\ n\ }Y_1\xrightarrow{\ 2\ }Y_2\xrightarrow{\ n\ }E_K\) through cyclic
          isogenies, the first two defined over \(K\) and the last one defined over \(K\) up to sign. Thus
          \(Y_2\) is \(K\)-isogenous to a quadratic twist of \(E_K\).
\item As in Step 1, \(\tau(\mu)=\pm\hat\mu\), and comparing the unique cyclic factorizations of
          \(\tau(\mu)\) and \(\hat\mu\) gives \(Y_1^\tau\cong Y_2\) over \(\bar\Q\).
\item Put \(X_0=Y_2\), a curve over \(K\). Let \(\varepsilon:G_K\to\{\pm1\}\) be the character by which \(G_K\) acts
          on the last factor \(Y_2\to E_K\) of \(\mu\); this factor is then a \(K\)-isogeny \(Y_2\to E_K\otimes\varepsilon\),
          so \(A_K\sim_K X_0\otimes\Z[i](\varepsilon\psi_K)\), and \(\varepsilon\psi_K\) has order \(4\) because
          \((\varepsilon\psi_K)^2=\psi_K^2\neq1\). We replace \(\psi_K\) by \(\varepsilon\psi_K\), so that
          \(A_K\sim_K X_0\otimes\Z[i](\psi_K)\) with \(\psi_K\) of order \(4\). Let \(\nu:X_0\to Y_1\) be the dual of \(Y_1\to Y_2\), a \(K\)-rational cyclic isogeny of
          degree \(2\). Then \(Y_1\cong X_0^\tau\) over \(\bar\Q\) (as \(\tau^2\in G_K\)); both curves are defined
          over \(K\) and have \(j\neq0,1728\), so \(Y_1\cong_K(X_0^\tau)^\eta\) for a character
          \(\eta:G_K\to\{\pm1\}\).
\end{enumerate}

\textbf{Step 3: the \(13\)-structure.} Let \(\iota\in\Z_{13}\) with \(\iota^2=-1\), and
      \(\psi_{K,\pm}=(i\mapsto\pm\iota)\circ\psi_K\). Then
      \(V_{13}A_K\cong V\otimes\psi_{K,+}\oplus V\otimes\psi_{K,-}\) with \(V=V_{13}X_0\). Write \(\bar\chi=\bar\chi_{13}\).

\begin{enumerate}
\item A rational point of order \(13\) gives a trivial constituent, so by (F4) \(X_0[13]^{\mathrm{ss}}\) contains
          \(\alpha=\bar\psi_{K,+}^{-1}\) or \(\bar\psi_{K,-}^{-1}\), a character of order \(4\) (reduction is
          injective on \(\mu_4\)). Hence \(X_0[13]\) is reducible over \(K\); let \(C\subset X_0[13]\) be a
          \(K\)-rational line, with character \(\beta\). Then \(X_0[13]^{\mathrm{ss}}=\beta\oplus\bar\chi\beta^{-1}\)
          and \(\alpha\in\{\beta,\bar\chi\beta^{-1}\}\).
\item As \(G_K\)-modules, \(X_0^\tau[13]\cong X_0[13]^\tau\); and via \(\nu\), which has degree prime to \(13\),
          \(X_0^\tau[13]\cong Y_1[13]\otimes\eta\) \({}\cong X_0[13]\otimes\eta\). Comparing semisimplifications,
          \(\beta^\tau\in\{\beta\eta,\ \bar\chi\beta^{-1}\eta\}\).
\item The second option is impossible: it gives \(\beta\beta^\tau=\bar\chi\eta\). If \(\beta=\alpha\), the left
          side has order dividing \(4\); if \(\beta=\bar\chi\alpha^{-1}\), it gives \(\bar\chi=\alpha\alpha^\tau\eta\).
          Either way \(\bar\chi^4|_{G_K}=1\), whereas \(\bar\chi|_{G_K}\) has order \(12\) or \(6\).
\item Hence \(\beta^\tau=\beta\eta\). Also \(\beta\neq\bar\chi\beta^{-1}\), since \(\bar\chi=\beta^2\) would
          again force \(\bar\chi^4|_{G_K}=1\). So \(\tau(C)\) is the unique \(G_K\)-stable line of \(X_0^\tau[13]\)
          with character \(\beta^\tau\), and under \(Y_1\cong(X_0^\tau)^\eta\) it corresponds to \(\nu(C)\), which
          has character \(\beta\eta=\beta^\tau\) there.
\end{enumerate}

\textbf{Step 4: a rational point of \(X_0(26)/w_2\).} Put \(x=(X_0,\ C+\ker\nu)\in Y_0(26)(K)\). Then

\[
w_2(x)=\bigl(X_0/\ker\nu,\ \nu(C)+\nu(X_0[2])\bigr)\cong\bigl(Y_1,\ \nu(C)+\ker\hat\nu\bigr).
\]

Under \(Y_1\cong X_0^\tau\), \(\nu(C)\) corresponds to \(\tau(C)\) by Step 3. The subgroup \(\ker\hat\nu\)
      has quotient \(Y_1/\ker\hat\nu\cong X_0\), and \((\ker\nu)^\tau\subset X_0^\tau[2]\) has quotient
      \(Y_1^\tau\cong X_0^{\tau^2}=X_0\); by \hyperref[lemma-two-subgroup]{Lemma~10} they correspond. Hence
      \(\tau(x)=w_2(x)\) in \(X_0(26)\), and the image of \(x\) in \(X_0(26)/w_2\) is fixed by \(\Gal(K/\Q)\),
      i.e. rational.

\textbf{Computation.} By \href{https://genus-2-torsion.github.io/data/knowledge/sources/qsimple_13_x0_26.m}{\texttt{qsimple\_\allowbreak{}13\_\allowbreak{}x0\_\allowbreak{}26.m}}, Magma's model of \(X_0(26)\) and its involution
      \(w_2\) are checked against the Hecke traces and the modular polynomial \(\Phi_2\); the quotient
      \(X_0(26)/w_2\) is the elliptic curve 26a1 \(=[1,0,1,-5,-8]\), of rank \(0\) with torsion \(\Z/3\). Above
      its three rational points lie the four rational cusps of \(X_0(26)\) and one pair of conjugate points
      over \(\Q(\sqrt{13})\) whose \(j\)-invariant is a root of the Hilbert class polynomial \(H_{-52}\), so they
      are CM points. The point \(x\) is neither a cusp nor a CM point, because \(X_0\sim E\) has no CM. This
      contradiction shows that \textbf{\(\Gamma=C_4\) is impossible when \(A(\Q)[13]\neq0\)}. \qed

\section{The case \(\Gamma=C_3\): no principal polarization}\label{c3}

In this section \(A\) is of type (S), \(\Gamma=\langle r\rangle\cong C_3\), and \(A[13]^{\mathrm{ss}}\) has a
      trivial constituent. This hypothesis and the conclusion are invariant under \(\Q\)-isogeny.

\textbf{Structure.} The trace-one representative \(R\) of \(r\) satisfies \(R^2-R+1=0\) and \(R^3=-1\).
      Put \(w=-R\); then \(w^3=1\), \(\mathrm{Ad}(w)=\mathrm{Ad}(R)\), and \(w\in D^\Gamma=\End^0_{\Q}(A)\).
      Replacing \(A\) by the image of \(a\mapsto(Na,Nwa,Nw^2a)\) in \(A^3\) for suitable \(N\), on which \(w\)
      acts as the cyclic shift, we may assume \(\Z[w]\subset\End_{\Q}(A)\). Let \(k:G_{\Q}\to\Z/3\) with
      \(g|_L=r^{k(g)}\), and twist \(A\) by the homomorphism \(\xi(g)=w^{-k(g)}\). On the twist \(B\), \(g\) acts on
      \(D\) by \(\mathrm{Ad}(w^{-k}R^k)=\mathrm{Ad}((-1)^k)=\mathrm{id}\), so \(\End^0_{\Q}(B)=M_2(\Q)\) and
      \(B\sim_{\Q}E_1^2\) for an elliptic curve \(E_1/\Q\) without CM. The image of \(w\) in \(M_2(\Q)\) has order
      \(3\), so it is conjugate to multiplication by \(\zeta_3\) on \(\Z[\zeta_3]\cong\Z^2\). Untwisting,

\[
A\sim_{\Q}E_1\otimes\Z[\zeta_3](\psi),\qquad \psi(g)=\zeta_3^{k(g)},\quad \ker\psi=G_L.
\]

\textbf{\(E_1\) has a rational \(13\)-isogeny and no rational \(3\)-isogeny.}

\begin{enumerate}
\item Since \(13\) splits in \(\Q(\zeta_3)\), \(V_{13}A=V\otimes\psi_1\oplus V\otimes\psi_2\) with
          \(V=V_{13}E_1\) and characters \(\psi_i:G_{\Q}\to\mu_3(\Z_{13})\). A trivial constituent of
          \(A[13]^{\mathrm{ss}}\) forces \(E_1[13]^{\mathrm{ss}}\) to contain some \(\bar\psi_i^{-1}\), so \(E_1[13]\) is
          reducible and \(E_1\) has a rational \(13\)-isogeny.
\item A rational \(3\)-isogeny would then give a rational cyclic \(39\)-isogeny, excluded by (F9). So
          \(\bar\rho=\bar\rho_{E_1,3}\) is irreducible. It is absolutely irreducible: otherwise its image would
          lie in a non-split Cartan subgroup \(\F_9^\times\), whose only element of order \(2\) is \(-1\), of
          determinant \(1\), while complex conjugation has order \(2\) and determinant \(-1\).
\item \emph{Clifford.} \(\bar\rho|_{G_L}\) is semisimple, since \(G_L\) is normal of index \(3\). Suppose that
          over \(\bar\F_3\) it is \(\chi_1\oplus\chi_2\). If \(\chi_1\neq\chi_2\), the group \(G_{\Q}/G_L\) of order
          \(3\) fixes both isotypic lines and \(\bar\rho\) is reducible over \(\bar\F_3\). If \(\chi_1=\chi_2\),
          \(\bar\rho(G_L)\) is scalar and the projective image is generated by the class of one \(\bar\rho(g)\),
          whose eigenvalue ratio is a cube root of unity in characteristic \(3\), hence \(1\); so
          \(\bar\rho(g)\) is a scalar times a unipotent, and \(\bar\rho\) is either reducible or has scalar
          image, which complex conjugation excludes. Hence \textbf{\(\bar\rho_{E_1,3}|_{G_L}\) is absolutely
          irreducible}.
\end{enumerate}

\textbf{The \(3\)-adic obstruction.} Let \(\mathcal O=\Z_3[\zeta_3]\), a discrete valuation ring with
      uniformizer \(\pi=1-\zeta_3\), and \(T=T_3E_1\).

\begin{enumerate}
\item On \(T_3(E_1\otimes\Z[\zeta_3](\psi))=T\otimes\mathcal O\), \(g\) acts by \(\rho(g)\otimes\psi(g)\). For
          \(g\in G_L\) this is \(\rho(g)\otimes1\), and by (F5) these elements span \(M_2(\Z_3)\otimes1\). For
          \(\tau\) with \(\psi(\tau)=\zeta_3\), \((\rho(\tau)^{-1}\otimes1)(\rho(\tau)\otimes\zeta_3)=1\otimes\zeta_3\).
          So the \(\Z_3\)-algebra generated by \(G_{\Q}\) is \(M_2(\mathcal O)\).
\item By (F6) and (F1), every \(B'\sim_{\Q}A\) has \(T_3B'=\pi^k(T\otimes\mathcal O)\) inside \(V_3A\), for some
          \(k\in\Z\).
\item By (F3), \(\NS(A_{\bar\Q})\otimes\Q\) is the space of symmetric bilinear forms on
          \(\Q(\zeta_3)\cong\Q^2\), with \(g\) acting through multiplication by \(\psi(g)\), a rotation of order
          dividing \(3\). The symmetric square of the two-dimensional rotation representation is the trivial
          representation plus a two-dimensional one, so the invariants are spanned by
          \(h_0(a,b)=\operatorname{Tr}_{\Q(\zeta_3)/\Q}(a\bar b)\). By (F2), any polarization of any
          \(B'\sim_{\Q}A\) has Riemann form \(c\cdot e\otimes h_0\) on \(V_3A\), with \(c\in\Q^\times\) and \(e\) the
          Weil pairing of \(E_1\).
\item On \(\mathcal O\) with basis \(1,\zeta_3\), the form \(h_0\) has Gram matrix
          \(\left(\begin{smallmatrix}2&-1\\-1&2\end{smallmatrix}\right)\), of determinant \(3\); on
          \(\pi^k\mathcal O\), \(h_0(\pi^ka,\pi^kb)=3^kh_0(a,b)\), so \(\det(c\,h_0|_{\pi^k\mathcal O})=3^{2k+1}c^2\)
          has odd \(3\)-adic valuation. Since \(e\) is perfect on \(T\),
          \(\det(e\otimes ch_0)\) on \(T\otimes\pi^k\mathcal O\) is a unit times \((3^{2k+1}c^2)^2\), of
          valuation \(4k+2+4v_3(c)\neq0\). So the Riemann form is never perfect on \(T_3B'\).
\end{enumerate}

By (F2), \textbf{no abelian surface \(\Q\)-isogenous to \(A\) has a principal polarization}; this is
      Proposition 2(a). \qed

\section{The case \(\Gamma=C_2\): only product polarizations}\label{c2}

In this section \(A\) is of type (S), \(\Gamma=\Gal(K/\Q)=\langle\sigma\rangle\) with \(K=L\) quadratic,
      and \(A(\Q)[13]\neq0\). Write \(\bar\chi=\bar\chi_{13}\) and let \(\varepsilon_K\) be the quadratic
      character of \(K\).

\subsection{Structure}\label{c2-structure}

\begin{enumerate}
\item \emph{A \(13\)-point over \(K\).} Since \(\End^0(A_K)=D\cong M_2(\Q)\), \(A_K\sim_K E_1^2\) for an elliptic
          curve \(E_1\subset A_K\) over \(K\), and \(A_K/E_1\sim_K E_1\). By \hyperref[lemma-sequence]{Lemma~1}, some elliptic
          curve \(E'\sim_K E_1\) has a \(K\)-point \(P'\) of order \(13\).
\item \emph{A model over \(\Q\).} By Bosman--Bruin--Dujella--Najman~\cite[Theorem 11(1)]{BBDN}, \(E'\cong_K E'^\sigma\). So \(j(E')\in\Q\),
          and \(j(E')\neq0,1728\) since \(E'\) has no CM. Hence there are an elliptic curve \(E_0/\Q\) without
          CM and a character \(\kappa:G_K\to\{\pm1\}\) with \(E'\cong_K E_0\otimes\kappa\); and
          \(A_{\bar\Q}\sim E_0^2\).
\item \emph{The rational \(13\)-isogeny.} Let \(C_0\subset E_0[13]\) correspond to \(\langle P'\rangle\); then
          \(G_K\) acts on \(C_0\) by \(\kappa\). For \(\tau\in G_{\Q}\setminus G_K\), \(\tau(C_0)\) is \(G_K\)-stable with
          character \(\kappa^\tau\). If \(\tau(C_0)\neq C_0\), then \(E_0[13]\cong C_0\oplus\tau(C_0)\) as
          \(G_K\)-modules, and \(\bar\chi|_{G_K}=\kappa\kappa^\tau\) would have order at most \(2\); but it has order
          \(12\) or \(6\). So \(C_0\) is a rational subgroup of order \(13\), and by (F9) it is the only one.
\item \emph{Its character.} Let \(\varphi:G_{\Q}\to\F_{13}^\times\) be the character of \(C_0\). Then
          \(\varphi|_{G_K}=\kappa\), so \(\varphi^2\in\{\mathbf 1,\varepsilon_K\}\). If \(\varphi^2=\mathbf 1\), the
          quadratic twist \(E_0\otimes\varphi\) would have a rational point of order \(13\), contrary to (F9).
          \textbf{Hence \(\varphi\) has order \(4\) and \(\varphi^2=\varepsilon_K\).} Moreover \(E_0[13]^{\mathrm{ss}}=
          \varphi\oplus\bar\chi\varphi^{-1}\).
\item \emph{\(K\) is real.} For complex conjugation \(c\), \(\bar\rho_{E_0,13}(c)\) has eigenvalues \(1\) and
          \(-1\), and \(\varphi(c)\) is one of them; so \(\varepsilon_K(c)=\varphi(c)^2=1\). (This agrees with
          Bosman--Bruin--Dujella--Najman~\cite[Theorem 9(1)]{BBDN}, which is not needed here.)
\item \emph{The isogeny class.} By (F7) over \(\Q\), \(A\sim_{\Q}E_0\otimes\rho\) with projective image of
          \(\rho\) equal to \(\Gamma\cong C_2\). As in Section~\ref{c4}, Step 1, simplicity of \(A\) over \(\Q\) forces
          \(\rho(G_{\Q})=\langle g\rangle\) with \(g^2=-1\), i.e. \(A\sim_{\Q}E_0\otimes\Z[i](\psi')\) for a character
          \(\psi':G_{\Q}\to\mu_4\) of order \(4\). Fix \(\iota\in\Z_{13}\) with \(\iota^2=-1\) and \(\iota\equiv5\pmod{13}\),
          and for a character \(\psi:G_{\Q}\to\mu_4\) put \(\psi_\pm=(i\mapsto\pm\iota)\circ\psi\), so that
          \(\psi_-=\psi_+^{-1}\) and

\[
V_{13}\bigl(E_0\otimes\Z[i](\psi)\bigr)=V\otimes\psi_+\ \oplus\ V\otimes\psi_-,\qquad V=V_{13}E_0.
\]

Let \(\omega:\mu_4(\F_{13})\to\mu_4(\Z[i])\) be the isomorphism with \(\omega(5)=i\); then
          \(\bar\psi_+=\omega^{-1}\circ\psi\). A trivial residual constituent of \(V\otimes\psi'_\pm\) requires
          \(\bar\psi'_\pm\in\{\varphi^{-1},\ \bar\chi^{-1}\varphi\}\). Only \(\varphi^{-1}\) has order \(4\), because
          \((\bar\chi^{-1}\varphi)^4=\bar\chi^{-4}\neq\mathbf 1\). Hence \(\psi'=\psi^{\pm1}\) with \(\psi:=\omega\circ\varphi\),
          and since complex conjugation on \(\Z[i]\) identifies the twists by \(\psi\) and \(\psi^{-1}\),

\[
A\sim_{\Q}A_\psi:=E_0\otimes\Z[i](\psi),\qquad \psi=\omega\circ\varphi,\quad \psi^2=\varepsilon_K,\quad \bar\psi_+=\varphi,\ \bar\psi_-=\varphi^{-1}.
\]
\item \emph{Polarizations of \(A_\psi\).} By (F3), \(\NS((A_\psi)_{\bar\Q})=\Sym_2(\Z)\), and \(g\) acts by
          \(S\mapsto\psi(g)^tS\psi(g)\), where \(\psi(g)\) is a rotation by a multiple of \(90^\circ\) on
          \(\Z[i]=\Z\oplus\Z i\). The invariants are \(\Z\cdot I\): the \(G_{\Q}\)-invariant part of
          \(\NS((A_\psi)_{\bar\Q})\) is \(\Z\theta_0\), where \(\theta_0\) is the product polarization, which is
          principal. The same holds for \(A'_\psi:=(E_0/C_0)\otimes\Z[i](\psi)\), with \(\theta'_0\). By (F2), for every
          \(B\sim_{\Q}A_\psi\) the \(G_{\Q}\)-invariant part of \(\NS(B_{\bar\Q})\otimes\Q\) is one-dimensional.
\end{enumerate}

\subsection{Admissible primes \(\ell\neq13\)}\label{c2-lattices}

Call a prime \(\ell\neq13\) \emph{admissible} if \(\bar\rho_{E_0,\ell}|_{G_K}\) is absolutely irreducible.

\begin{statement}{Lemma 11.}\label{lemma-lattice}
Let \(\ell\neq13\) be admissible and \(\mathcal O=\Z_\ell[i]\). Every \(G_{\Q}\)-stable lattice in
          \(V_\ell A_\psi\) is \(\mathfrak a\,(T_\ell E_0\otimes\mathcal O)\) for a fractional \(\mathcal O\)-ideal
          \(\mathfrak a\).
\end{statement}

\emph{Proof.} For \(g\in G_K\), \(\psi(g)=\pm1\), so \(g\) acts by \((\psi(g)\rho(g))\otimes1\); these elements
      form a quadratic twist of \(\rho(G_K)\), whose reduction is still absolutely irreducible, and by (F5)
      they span \(M_2(\Z_\ell)\otimes1\). For \(\tau\notin G_K\), \(\psi(\tau)=\pm i\) and
      \((\rho(\tau)^{-1}\otimes1)(\rho(\tau)\otimes\psi(\tau))=\pm1\otimes i\). So the \(\Z_\ell\)-algebra generated by
      \(G_{\Q}\) is \(M_2(\mathcal O)\), and (F6) applies; \(\mathcal O\) is \(\Z_\ell\times\Z_\ell\), the unramified
      quadratic extension of \(\Z_\ell\), or \(\Z_2[i]\). \qed

\subsection{The prime \(13\)}\label{c2-thirteen}

Let \(e_\pm\) be the idempotents of \(\Z_{13}[i]=\Z_{13}e_+\oplus\Z_{13}e_-\) with \(ie_\pm=\pm\iota e_\pm\), so
      that \(V_{13}A_\psi=V\otimes\psi_+e_+\oplus V\otimes\psi_-e_-\). Since \(E_0[13]^{\mathrm{ss}}=\varphi\oplus\bar\chi\varphi^{-1}\)
      and \(\varphi^2=\varepsilon_K\):

\begin{itemize}
\item the residual constituents of \(V\otimes\psi_+\) are \(\{\varepsilon_K,\ \bar\chi\}\);
\item those of \(V\otimes\psi_-\) are \(\{\mathbf 1,\ \bar\chi\varepsilon_K\}\).
\end{itemize}

These sets are disjoint, because \(\varepsilon_K\neq\mathbf 1\) and \(\bar\chi\notin\{\mathbf 1,\varepsilon_K\}\).

\begin{statement}{Lemma 12 (splitting).}\label{lemma-split}
Let \(X,Y\) be finite-dimensional \(\Q_\ell[G]\)-modules whose residual constituents are disjoint.
          Then every \(G\)-stable lattice \(\Lambda\subset X\oplus Y\) equals \((\Lambda\cap X)\oplus(\Lambda\cap Y)\).
\end{statement}

\emph{Proof.} Let \(p_X,p_Y\) be the projections. The map \(\lambda\mapsto p_X(\lambda)\) induces
      isomorphisms of \(G\)-modules

\[
\Lambda/\bigl((\Lambda\cap X)\oplus(\Lambda\cap Y)\bigr)\cong p_X\Lambda/(\Lambda\cap X)\cong p_Y\Lambda/(\Lambda\cap Y).
\]

By (F4), the constituents of the middle module lie among those of \(X\), and those of the right-hand
      module among those of \(Y\). Hence all three are zero. \qed

Now let \((B,\lambda)\) be principally polarized with \(B\sim_{\Q}A_\psi\), and fix a \(\Q\)-isogeny
      \(f:B\to A_\psi\); we identify \(T_\ell B\) with \(f_*T_\ell B\subset V_\ell A_\psi\).

\begin{enumerate}
\item By Lemma 12, \(T_{13}B=\Lambda_+e_+\oplus\Lambda_-e_-\) with \(G_{\Q}\)-stable lattices
          \(\Lambda_\pm\subset V\); twisting by the scalar characters \(\psi_\pm\) does not change stability.
\item The \(G_{\Q}\)-stable lattices in \(V\) correspond, up to homothety, to rational cyclic isogenies of
          \(13\)-power degree from \(E_0\). By (F9) the only nontrivial one is \(E_0\to E_0/C_0\). So each
          \(\Lambda_\pm\) is homothetic to \(T=T_{13}E_0\) or to \(T'\), where \(T'\supset T\) and \(T'/T=C_0\);
          that is, \(T'=T_{13}(E_0/C_0)\) under the isogeny.
\item Let \(E_\theta\) be the Riemann form of \(\theta_0\); it is \(e\otimes h_0\) with
          \(h_0(z,z')=\operatorname{Re}(z\bar z')=\tfrac12\operatorname{Tr}(z\bar z')\). Complex conjugation on
          \(\Z[i]\) exchanges \(e_+\) and \(e_-\), so \(h_0(e_\pm,e_\pm)=0\) and \(h_0(e_+,e_-)=\tfrac12\), and

\[
E_\theta(x_+e_++x_-e_-,\,y_+e_++y_-e_-)=\tfrac12\bigl(e(x_+,y_-)+e(x_-,y_+)\bigr).
\]
\item By Section~\ref{c2-structure}, item 7, the Riemann form of \(\lambda\) is \(cE_\theta\) for some
          \(c\in\Q^\times\). It is perfect on \(T_{13}B\) if and only if \(\Lambda_-=(2/c)\Lambda_+^*\), where
          \(\Lambda^*=\{v\in V:e(v,\Lambda)\subseteq\Z_{13}\}\).
\item \(T^*=T\). If \(T=\langle c_1,d_1\rangle\) with \(C_0=\langle\bar c_1\rangle\), then
          \(T'=\langle c_1/13,d_1\rangle\) and \(T'^*=\langle c_1,13d_1\rangle=13T'\).
\item Hence \((\Lambda_+,\Lambda_-)=(13^aT,13^bT)\) or \((13^aT',13^bT')\): the mixed shapes are not
          dual to each other. In the first case \(T_{13}B=\mathfrak a_{13}(T\otimes\Z_{13}[i])\) \({}=\mathfrak a_{13}T_{13}A_\psi\)
          with \(\mathfrak a_{13}=13^a\Z_{13}e_+\oplus13^b\Z_{13}e_-\). In the second case the same holds with
          \(T'\), i.e. relative to \(A'_\psi\), where we identify \(V_\ell A'_\psi=V_\ell A_\psi\) via
          \(E_0\to E_0/C_0\); for \(\ell\neq13\) the lattices \(T_\ell A'_\psi\) and \(T_\ell A_\psi\) agree.
\end{enumerate}

\subsection{Conclusion, assuming every prime is admissible}\label{c2-conclusion}

\begin{statement}{Proposition 13.}\label{prop-c2}
If every prime \(\ell\neq13\) is admissible, then every principally polarized \((B,\lambda)\) with
          \(B\sim_{\Q}A_\psi\) is isomorphic to \((A_\psi,\theta_0)\) or to \((A'_\psi,\theta'_0)\). In particular
          \((B_{\bar\Q},\lambda)\) is a product of elliptic curves with the product polarization, and
          \((B,\lambda)\) is not a Jacobian.
\end{statement}

\emph{Proof.} By Lemmas 11 and 12 and the discussion above, \(T_\ell B=\mathfrak a_\ell\,T_\ell A^{(\prime)}_\psi\)
      for every \(\ell\), with \(\mathfrak a_\ell=\Z_\ell[i]\) for almost all \(\ell\). These local ideals come from a
      fractional ideal \(\mathfrak a\) of the principal ideal domain \(\Z[i]\), say \(\mathfrak a=\alpha\Z[i]\). Since
      \(\alpha\in\Q(i)\subseteq\End^0_{\Q}(A_\psi)\), (F1) gives \(B\cong A^{(\prime)}_\psi\). The only principal
      polarization on \(A^{(\prime)}_\psi\) is \(\theta_0^{(\prime)}\), since \(n\theta_0^{(\prime)}\) has degree \(n^4\).
      Geometrically \((A_\psi,\theta_0)\cong(E_0^2,\text{product})\) and \((A'_\psi,\theta'_0)\cong((E_0/C_0)^2,\text{product})\);
      by (F8) these are not Jacobians. \qed

\subsection{Every prime \(\ell\neq13\) is admissible}\label{c2-admissible}

Suppose that \(\bar\rho=\bar\rho_{E_0,\ell}|_{G_K}\) is not absolutely irreducible. Then either it is
      reducible over \(\F_\ell\), or it is irreducible but its commutant contains \(\F_{\ell^2}\), so that its
      image lies in a non-split Cartan subgroup \(C_{ns}(\ell)\cong\F_{\ell^2}^\times\). For \(\ell=2\),
      \(\GL_2(\F_2)\cong S_3\): the image \(S_3\) is absolutely irreducible, images of order at most \(2\) fix a
      nonzero vector (reducible case), and the remaining case is the image \(C_3=C_{ns}(2)\).

\textbf{(a) Reducible over \(\F_\ell\).} A \(G_K\)-stable line \(L_\ell\subset E_0[\ell]\) together with \(C_0\)
      gives a \(K\)-rational cyclic \(13\ell\)-isogeny of \(E_0\), a curve without CM with \(j(E_0)\in\Q\).

\begin{itemize}
\item For \(\ell=2\), the cubic \(2\)-division polynomial of \(E_0\) has a root in the quadratic field \(K\),
          hence a rational root; so \(E_0\) has a rational cyclic \(26\)-isogeny, contrary to (F9).
\item For odd \(\ell\neq13\), Vukorepa~\cite[Theorem 1]{Vuk} states: if \(E\) is an elliptic curve without CM with
          \(j(E)\in\Q\) and a cyclic \(n\)-isogeny defined over a quadratic field, then \(n\le18\) with \(n\neq14\), or \(n\) is one of
          \(20\), \(21\), \(24\), \(25\), \(32\), \(36\), \(37\). Since \(13\ell\ge39\) is not in this list, this case is impossible.
          In Vukorepa~\cite{Vuk} the relevant cases are treated in Lemma 4.1 (``Case \(q=13\)'', with \(q=17,37\) for
          \(\ell=17,37\)): for \(\ell=5\) through the quadratic points of \(X_0(65)\) determined in
          Box~\cite[§4.5]{Box}; for \(\ell=7\) through the quadratic points of \(X_0(91)\) determined in
          \cite[§13]{Vuk}; for \(\ell=11,17,37\) because the isogeny descends to \(\Q\); and for \(\ell=3\)
          through a genus-\(2\) curve of rank \(0\) on which \(j\) is a cube (part (4) of
          \href{https://genus-2-torsion.github.io/data/knowledge/sources/qsimple_13_x0_13.m}{\texttt{qsimple\_\allowbreak{}13\_\allowbreak{}x0\_\allowbreak{}13.m}} re-checks this curve). For the remaining primes there is no
          \(K\)-rational \(\ell\)-isogeny at all (\cite[§3]{Vuk}).
\end{itemize}

\textbf{(b) \(\ell\) odd, image in \(C_{ns}(\ell)\).} Complex conjugation \(c\) lies in \(G_K\) because \(K\) is
      real. Then \(\bar\rho(c)\in\F_{\ell^2}^\times\) has order dividing \(2\), so \(\bar\rho(c)=\pm1\), of determinant
      \(1\); but \(\det\bar\rho(c)=\bar\chi_\ell(c)=-1\).

\textbf{(c) \(\ell=2\), image \(C_3\).} Then \(G_K\) permutes the roots of the \(2\)-division cubic by even
      permutations, so the discriminant \(\Delta(E_0)\) is a square in \(K\): \(\Delta(E_0)\in\Q^{\times2}\) or
      \(\Delta(E_0)\in d_K\Q^{\times2}\), where \(K=\Q(\sqrt{d_K})\). Since \(j-1728=c_6^2/\Delta\) and
      \(j(E_0)\neq1728\), \(\Delta(E_0)\equiv j(E_0)-1728\) modulo \(\Q^{\times2}\).

\begin{itemize}
\item \textbf{(c1) \(\Delta(E_0)\in d_K\Q^{\times2}\).} Let \(C\) be the genus-\(2\) curve
          \(y^2+(x^3+x^2+1)y=x^2+x\). By \href{https://genus-2-torsion.github.io/data/knowledge/sources/qsimple_13_x1_13.m}{\texttt{qsimple\_\allowbreak{}13\_\allowbreak{}x1\_\allowbreak{}13.m}}, the Tate normal form
          \(Y^2+(1-c)XY-bY=X^3-bX^2\) with \(r=1-xy\), \(s=1-xy/(y+1)\), \(b=rs(r-1)\), \(c=s(r-1)\) has \((0,0)\)
          of exact order \(13\) over the function field of \(C\). The induced map \(C\to X_1(13)\) is
          non-constant between curves of genus \(2\), hence an isomorphism over \(\Q\), and the \(j\)-invariant
          is a rational function \(j(x)\in\Q(x)\) with denominator \(x^{13}(x+1)^{13}(x^3+4x^2+x-1)\). The
          hyperelliptic involution is \((x,y)\mapsto(x,-y-x^3-x^2-1)\), with quotient map \(x\).

The pair \((E',\pm P')\) is a point of \(Y_1(13)\) with residue field \(K\), since
          \(Y_1(13)(\Q)=\emptyset\) by (F9). By Bosman--Bruin--Dujella--Najman~\cite[Lemmas 12 and 13]{BBDN} (\(J_1(13)(\Q)\cong\Z/19\Z\), every
          class being represented by a divisor supported on the cusps), as in the proof of
          \cite[Theorem 9(1)]{BBDN}, this point lies in a fibre of \(x\) over a rational point. The
          two points at infinity of \(C\) are rational, so \(x=t\in\Q\), \(K=\Q(\sqrt{D(t)})\) with
          \(D(t)=(t^3+t^2+1)^2+4(t^2+t)\), and \(j(E_0)=j(E')=j(t)\).

The script also verifies

\[
j(t)-1728=-\frac{D(t)\,g(t)^2}{t^{13}(t+1)^{13}(t^3+4t^2+t-1)}
\]

with \(g\in\Q[t]\) of degree \(18\); here \(g(t)\neq0\) because \(j\neq1728\). So the condition
          \((j(t)-1728)D(t)\in\Q^{\times2}\) means that \(t\) is the \(x\)-coordinate of a rational point on

\[
C_1:\ s^2=-t(t+1)(t^3+4t^2+t-1)=-t^5-5t^4-5t^3+t.
\]

The Mordell--Weil rank of \(J(C_1)(\Q)\) is \(0\) unconditionally: the cubic field
          \(\Q[t]/(t^3+4t^2+t-1)\) of discriminant \(169\) has trivial class group (computed with
          \texttt{Proof := {\char34}Full{\char34}}), and \(\#\mathrm{Sel}_2(J(C_1))=\#J(C_1)(\Q)[2]=4\). Chabauty's method for rank
          \(0\) gives \(C_1(\Q)=\{(-1,0),(0,0),\infty\}\), i.e. \(t\in\{-1,0,\infty\}\), which are cusps (poles of
          \(j\)). Contradiction.
\item \textbf{(c2) \(\Delta(E_0)\in\Q^{\times2}\).} Since \(E_0\) has a rational \(13\)-isogeny,
          \(j(E_0)=j_{13}(u)\) for some \(u\in\Q^\times\), where

\[
j_{13}(u)=\frac{(u^2+5u+13)(u^4+7u^3+20u^2+19u+1)^3}{u}
\]

parametrizes \(X_0(13)\) (see \cite[§3.3]{Vuk}). The script \href{https://genus-2-torsion.github.io/data/knowledge/sources/qsimple_13_x0_13.m}{\texttt{qsimple\_\allowbreak{}13\_\allowbreak{}x0\_\allowbreak{}13.m}} checks
          that \(\Phi_{13}(j_{13}(u),j_{13}(13/u))\) vanishes identically and that \(\deg j_{13}=14\), so
          \(u\mapsto(j_{13}(u),j_{13}(13/u))\) is a birational parametrization of \(X_0(13)\). It also verifies
          \(j_{13}(u)-1728\) \({}=(u^2+6u+13)h(u)^2/u\) with \(h\) of degree \(6\) and no rational root, so
          \(j-1728\in\Q^{\times2}\) means \(s^2=u^3+6u^2+13u\). With \(u=X-2\) this is \(Y^2=X^3+X-10\), the curve
          52a1, of rank \(0\) with torsion \(\Z/2\Z\). Its rational points are \(\infty\) and \((u,s)=(0,0)\), which
          are cusps. Contradiction.
\end{itemize}

Hence every prime \(\ell\neq13\) is admissible, and Proposition 13 proves Proposition 2(b). This
      completes the proof of Proposition 2 and of the Theorem. \qed

\section{Computations}\label{computations}

Each script below is self-contained, asserts every fact it reports, and ends with
      \texttt{ALL\_CHECKS\_PASSED}. They were run in Magma V2.30-1 on 21 October 2026; the logs are
      linked next to the scripts. No result depends on the generalized Riemann hypothesis: Magma's
      class group computations are rigorous unless \texttt{SetClassGroupBounds({\char34}GRH{\char34})} is called,
      which none of the scripts does.

\begin{itemize}
\item \href{https://genus-2-torsion.github.io/data/knowledge/sources/qsimple_13_cm_weil.m}{\texttt{qsimple\_\allowbreak{}13\_\allowbreak{}cm\_\allowbreak{}weil.m}} (\href{https://genus-2-torsion.github.io/data/knowledge/sources/qsimple_13_cm_weil.log}{\texttt{qsimple\_\allowbreak{}13\_\allowbreak{}cm\_\allowbreak{}weil.log}}): the exact enumeration of the
          Weil polynomials of abelian surfaces over \(\F_2\), the polynomial \(T^4+2T^3+2T^2+4T+4\) and its root
          field \(\Q(\zeta_{12})\), and the facts about \(\Q(i)\), \(\Q(\sqrt{-3})\) and \(E_M\) used in
          Section~\ref{type-cm}.
\item \href{https://genus-2-torsion.github.io/data/knowledge/sources/qsimple_13_x0_39.m}{\texttt{qsimple\_\allowbreak{}13\_\allowbreak{}x0\_\allowbreak{}39.m}} (\href{https://genus-2-torsion.github.io/data/knowledge/sources/qsimple_13_x0_39.log}{\texttt{qsimple\_\allowbreak{}13\_\allowbreak{}x0\_\allowbreak{}39.log}}): the model of \(X_0(39)\), the moduli check of
          \(w_3\), \(X_0(39)/w_3\cong\) 39a1 of rank \(0\), and the fibres over its four rational points
          (Section~\ref{c6}); hence \(X_0(39)(\Q)\) consists of cusps (F9).
\item \href{https://genus-2-torsion.github.io/data/knowledge/sources/qsimple_13_x0_26.m}{\texttt{qsimple\_\allowbreak{}13\_\allowbreak{}x0\_\allowbreak{}26.m}} (\href{https://genus-2-torsion.github.io/data/knowledge/sources/qsimple_13_x0_26.log}{\texttt{qsimple\_\allowbreak{}13\_\allowbreak{}x0\_\allowbreak{}26.log}}): the model of \(X_0(26)\), the moduli check of
          \(w_2\), \(X_0(26)/w_2\cong\) 26a1 of rank \(0\), and the fibres over its three rational points,
          including \(H_{-52}\) (Section~\ref{c4}); hence \(X_0(26)(\Q)\) consists of cusps (F9).
\item \href{https://genus-2-torsion.github.io/data/knowledge/sources/qsimple_13_x1_13.m}{\texttt{qsimple\_\allowbreak{}13\_\allowbreak{}x1\_\allowbreak{}13.m}} (\href{https://genus-2-torsion.github.io/data/knowledge/sources/qsimple_13_x1_13.log}{\texttt{qsimple\_\allowbreak{}13\_\allowbreak{}x1\_\allowbreak{}13.log}}): the model of \(X_1(13)\), the \(j\)-map, the
          factorization of \(j-1728\), and the unconditional determination of \(C_1(\Q)\)
          (Section~\ref{c2-admissible}, (c1)); also the example of Section~\ref{remarks}.
\item \href{https://genus-2-torsion.github.io/data/knowledge/sources/qsimple_13_x0_13.m}{\texttt{qsimple\_\allowbreak{}13\_\allowbreak{}x0\_\allowbreak{}13.m}} (\href{https://genus-2-torsion.github.io/data/knowledge/sources/qsimple_13_x0_13.log}{\texttt{qsimple\_\allowbreak{}13\_\allowbreak{}x0\_\allowbreak{}13.log}}): the parametrization \(j_{13}\) of \(X_0(13)\), the
          factorization of \(j_{13}-1728\), the curve 52a1 (Section~\ref{c2-admissible}, (c2)), and the cross-check of
          the genus-\(2\) curve \(w^3=u^2(u^2+5u+13)\) used in Vukorepa~\cite{Vuk} for \(\ell=3\).
\end{itemize}

The inputs from the literature that were not recomputed here are
      Bosman--Bruin--Dujella--Najman~\cite[Theorem 11(1), Lemmas 12 and 13]{BBDN}, Vukorepa~\cite[Theorem 1]{Vuk} (including the computations
      on \(X_0(65)\) in Box~\cite[§4.5]{Box} and on \(X_0(91)\) in \cite[§13]{Vuk}), Mazur's and Kenku's
      theorems in (F9), Faltings' theorems (F10), Tate's theorem in Section~\ref{type-cm}, and the criterion of
      Néron--Ogg--Shafarevich.

\section{Consequence for the census}\label{census}

The group \href{https://genus-2-torsion.github.io/group.html?g=13}{\([13]\)} is impossible in the class ``geometrically split, simple over \(\Q\)''. It was
      already impossible in the class ``split over \(\Q\)'' by Mazur's theorem
      (\href{https://genus-2-torsion.github.io/proofs/qsplit-prime-obstruction.html}{prime torsion obstruction}). Hence no geometrically split
      genus-\(2\) Jacobian over \(\Q\) has a rational point of order \(13\), and the combined geometrically
      split \(\infty\)? cell of \href{https://genus-2-torsion.github.io/group.html?g=13}{\([13]\)} becomes \(\emptyset\).

The theorem applies to every group containing a point of order \(13\). Among the census groups these
      are \href{https://genus-2-torsion.github.io/group.html?g=13}{\([13]\)}, \href{https://genus-2-torsion.github.io/group.html?g=26}{\([26]\)}, \href{https://genus-2-torsion.github.io/group.html?g=2.26}{\([2,26]\)} and \href{https://genus-2-torsion.github.io/group.html?g=39}{\([39]\)}. The last three are also excluded, in both geometrically
      split classes, by \href{https://genus-2-torsion.github.io/proofs/geometric-26-39-obstruction.html}{the proof for orders 26 and 39}, which shows
      independently that no geometrically split abelian surface over \(\Q\), principally polarized or not,
      has a rational point of order \(26\) or \(39\). For \href{https://genus-2-torsion.github.io/group.html?g=13}{\([13]\)} the principal polarization is essential;
      see Section~\ref{remarks}.

The argument places no restriction on the geometric splitting degree, the field of definition of
      the geometric endomorphisms, or the integral endomorphism ring.

\section{Remarks}\label{remarks}

\begin{enumerate}
\item \emph{The Jacobian hypothesis cannot be dropped.} Let \(E\) be an elliptic curve without CM over a
          quadratic field \(K\) with a point of order \(13\). By Bosman--Bruin--Dujella--Najman~\cite[Theorem 9(1)]{BBDN}, \(K\) is real and
          \(E\) has false complex multiplication by \(\Q(\sqrt{-1})\), which by the definition in \cite[§4.1]{BBDN} means that
          \(\End^0_{\Q}(\operatorname{Res}_{K/\Q}E)\) contains \(\Q(i)\); this algebra has dimension \(2\), because
          \(E\) has no CM and \(E\cong_K E^\sigma\) \cite[Theorem 11(1)]{BBDN}, so
          \(\End^0_{\Q}(\operatorname{Res}_{K/\Q}E)\cong\Q(i)\). So it is simple over \(\Q\) and geometrically
          split, and its group of rational points \(E(K)\) contains a point of order \(13\). Such curves exist: the
          two points of the curve \(C\) of Section~\ref{c2-admissible} with \(x=-2\) are defined over \(K=\Q(\sqrt{17})\)
          and give such a curve \(E\), with \(j\)-invariant \(-60698457/40960\); this is the \(j\)-invariant of the curve
          \([1,-1,0,-139,965]\) over \(\Q\), which has a rational \(13\)-isogeny
          (\href{https://genus-2-torsion.github.io/data/knowledge/sources/qsimple_13_x1_13.m}{\texttt{qsimple\_\allowbreak{}13\_\allowbreak{}x1\_\allowbreak{}13.m}}, part (5)). These surfaces are of type \(\Gamma=C_2\); by Proposition 2(b), all
          principal polarizations in their isogeny classes are product polarizations.
\item \emph{What is proved for abelian surfaces.} For a geometrically split abelian surface \(A/\Q\) that
          is simple over \(\Q\) and has \(A(\Q)[13]\neq0\), with any polarization: \(A\) is not of type (N) or (CM),
          and in type (S) \(\Gamma\notin\{C_1,C_4,C_6,V_4,D_6,D_8,D_{12}\}\). If \(\Gamma=C_3\), no surface in the
          \(\Q\)-isogeny class of \(A\) is principally polarized; if \(\Gamma=C_2\), every principally polarized
          surface in it is geometrically a product with the product polarization.
\item \emph{Earlier ingredients.} The argument for \(\Gamma=C_6\) and for non-cyclic \(\Gamma\) is the
          \(13\)-part of the argument used for orders \(26\) and \(39\); the construction of the curve \(E_c\)
          in Section~\ref{c4} is a special case of Elkies's reduction of quadratic \(K\)-curves
          (Elkies~\cite{Elk}); the argument is given here in full.
\end{enumerate}
\begin{thebibliography}{Ken80}

\bibitem[BBDN]{BBDN}
J. Bosman, P. Bruin, A. Dujella and F. Najman,
\emph{Ranks of elliptic curves with prescribed torsion over number fields},
Int. Math. Res. Not. IMRN \textbf{2014}, no.~11, 2885--2923;
     \href{https://arxiv.org/abs/1201.0252}{arXiv:1201.0252v1}. Theorem~9(1), Theorem~11(1),
     Lemmas~12 and~13 and \S4.1 are cited in the numbering of arXiv:1201.0252v1.

\bibitem[Vuk]{Vuk}
B. Vukorepa,
\emph{Isogenies over quadratic fields of elliptic curves with rational \(j\)-invariant},
\href{https://arxiv.org/abs/2203.10672}{arXiv:2203.10672v1} (2022), Theorem~1; Lemma~4.1 and
     \S\S3, 13.

\bibitem[Box]{Box}
J. Box,
\emph{Quadratic points on modular curves with infinite Mordell--Weil group},
Math. Comp. \textbf{90} (2021), 321--343, \S4.5 (numbering of arXiv:1906.05206v3).
     \href{https://arxiv.org/abs/1906.05206}{arXiv:1906.05206}.

\bibitem[Maz]{Maz}
B. Mazur,
\emph{Modular curves and the Eisenstein ideal},
Publ. Math. IH\'ES \textbf{47} (1977), 33--186, Theorem~(7$'$).
     \href{https://www.numdam.org/item/PMIHES_1977__47__33_0/}{Numdam}.

\bibitem[Ken79]{Ken79}
M. A. Kenku,
\emph{The modular curve \(X_0\)(39) and rational isogeny},
Math. Proc. Cambridge Philos. Soc. \textbf{85} (1979), 21--23.

\bibitem[Ken80]{Ken80}
M. A. Kenku,
\emph{The modular curve \(X_0\)(169) and rational isogeny},
J. London Math. Soc. (2) \textbf{22} (1980), 239--244.

\bibitem[Fal]{Fal}
G. Faltings,
\emph{Endlichkeitss\"atze f\"ur abelsche Variet\"aten \"uber Zahlkörpern},
Invent. Math. \textbf{73} (1983), 349--366, \S5, Satz~3 and Korollar~1.

\bibitem[ST]{ST}
J.-P. Serre and J. Tate,
\emph{Good reduction of abelian varieties},
Ann. of Math. (2) \textbf{88} (1968), 492--517, \S1, Theorem~1 and Corollaries~2--3 (p.~493); \S2, Theorem~2(i) (p.~496).

\bibitem[Tat]{Tat}
J. Tate,
\emph{Endomorphisms of abelian varieties over finite fields},
Invent. Math. \textbf{2} (1966), 134--144, Theorem~2(c).

\bibitem[Elk]{Elk}
N. D. Elkies,
\emph{On elliptic \(K\)-curves},
in: \emph{Modular curves and abelian varieties}, Progr. Math. \textbf{224}, Birkh\"auser, Basel, 2004,
     81--91.

\end{thebibliography}

\end{document}
